% Ngân hàng đề Toán 12 (KNTT) -- Chương II, Bài 8, Dạng 09
% Dạng: Hai vectơ cùng phương, ba điểm thẳng hàng, hình bình hành
%% Phương án ĐẦU TIÊN trong \cauTN luôn là đáp án đúng (gói trondethi tự đảo).

\cauTN[id=C2-B8-D09-NB-001,muc=NB,loigiai={Ba điểm thẳng hàng khi hai vectơ tạo bởi chúng cùng phương. Ta có $\overrightarrow{AB}=(-9;-3;-3)$ và $\overrightarrow{AE}=(-6;-2;-2)$ cùng phương nên $A$, $B$, $E$ thẳng hàng. Các bộ còn lại không thoả mãn.}]{Trong không gian $Oxyz$, cho $A(8;6;5)$, $B(-1;3;2)$, $C(-3;2;-1)$, $D(-2;-2;4)$, $E(2;4;3)$. Ba điểm nào dưới đây thẳng hàng?}
  {$A, B, E$.}
  {$A, B, D$.}{$A, D, E$.}{$B, D, E$.}

\cauTN[id=C2-B8-D09-NB-002,muc=NB,loigiai={Ba điểm thẳng hàng khi hai vectơ tạo bởi chúng cùng phương. Ta có $\overrightarrow{BC}=(-6;-1;-2)$ và $\overrightarrow{BE}=(-12;-2;-4)$ cùng phương nên $B$, $C$, $E$ thẳng hàng. Các bộ còn lại không thoả mãn.}]{Trong không gian $Oxyz$, cho $A(1;-1;0)$, $B(4;-1;-2)$, $C(-2;-2;-4)$, $D(2;2;-4)$, $E(-8;-3;-6)$. Ba điểm nào dưới đây thẳng hàng?}
  {$B, C, E$.}
  {$A, C, E$.}{$A, B, E$.}{$A, C, D$.}

\cauTN[id=C2-B8-D09-NB-003,muc=NB,loigiai={Ba điểm thẳng hàng khi hai vectơ tạo bởi chúng cùng phương. Ta có $\overrightarrow{AB}=(9;-6;9)$ và $\overrightarrow{AE}=(6;-4;6)$ cùng phương nên $A$, $B$, $E$ thẳng hàng. Các bộ còn lại không thoả mãn.}]{Trong không gian $Oxyz$, cho $A(-5;8;-7)$, $B(4;2;2)$, $C(3;-1;0)$, $D(-4;3;-1)$, $E(1;4;-1)$. Ba điểm nào dưới đây thẳng hàng?}
  {$A, B, E$.}
  {$A, C, E$.}{$A, D, E$.}{$A, C, D$.}

\cauTN[id=C2-B8-D09-NB-004,muc=NB,loigiai={Ba điểm thẳng hàng khi hai vectơ tạo bởi chúng cùng phương. Ta có $\overrightarrow{BD}=(2;0;-1)$ và $\overrightarrow{BE}=(4;0;-2)$ cùng phương nên $B$, $D$, $E$ thẳng hàng. Các bộ còn lại không thoả mãn.}]{Trong không gian $Oxyz$, cho $A(3;0;0)$, $B(-4;0;3)$, $C(1;-3;3)$, $D(-2;0;2)$, $E(0;0;1)$. Ba điểm nào dưới đây thẳng hàng?}
  {$B, D, E$.}
  {$A, C, E$.}{$C, D, E$.}{$A, C, D$.}

\cauTN[id=C2-B8-D09-NB-005,muc=NB,loigiai={Ba điểm thẳng hàng khi hai vectơ tạo bởi chúng cùng phương. Ta có $\overrightarrow{BC}=(-2;-2;-2)$ và $\overrightarrow{BD}=(-4;-4;-4)$ cùng phương nên $B$, $C$, $D$ thẳng hàng. Các bộ còn lại không thoả mãn.}]{Trong không gian $Oxyz$, cho $A(-1;2;-1)$, $B(3;2;6)$, $C(1;0;4)$, $D(-1;-2;2)$, $E(4;1;3)$. Ba điểm nào dưới đây thẳng hàng?}
  {$B, C, D$.}
  {$A, B, C$.}{$B, C, E$.}{$A, C, D$.}

\cauTN[id=C2-B8-D09-NB-006,muc=NB,loigiai={Ba điểm thẳng hàng khi hai vectơ tạo bởi chúng cùng phương. Ta có $\overrightarrow{AB}=(2;2;2)$ và $\overrightarrow{AD}=(1;1;1)$ cùng phương nên $A$, $B$, $D$ thẳng hàng. Các bộ còn lại không thoả mãn.}]{Trong không gian $Oxyz$, cho $A(-1;-2;2)$, $B(1;0;4)$, $C(3;4;-3)$, $D(0;-1;3)$, $E(4;-3;-4)$. Ba điểm nào dưới đây thẳng hàng?}
  {$A, B, D$.}
  {$A, B, E$.}{$C, D, E$.}{$B, C, D$.}

\cauTN[id=C2-B8-D09-NB-007,muc=NB,loigiai={Ba điểm thẳng hàng khi hai vectơ tạo bởi chúng cùng phương. Ta có $\overrightarrow{AC}=(-3;3;-4)$ và $\overrightarrow{AE}=(3;-3;4)$ cùng phương nên $A$, $C$, $E$ thẳng hàng. Các bộ còn lại không thoả mãn.}]{Trong không gian $Oxyz$, cho $A(-2;0;-1)$, $B(1;-4;-2)$, $C(-5;3;-5)$, $D(2;-3;4)$, $E(1;-3;3)$. Ba điểm nào dưới đây thẳng hàng?}
  {$A, C, E$.}
  {$A, B, D$.}{$C, D, E$.}{$A, B, E$.}

\cauTN[id=C2-B8-D09-NB-008,muc=NB,loigiai={Ba điểm thẳng hàng khi hai vectơ tạo bởi chúng cùng phương. Ta có $\overrightarrow{AC}=(6;9;9)$ và $\overrightarrow{AD}=(2;3;3)$ cùng phương nên $A$, $C$, $D$ thẳng hàng. Các bộ còn lại không thoả mãn.}]{Trong không gian $Oxyz$, cho $A(1;-4;-1)$, $B(-3;3;-2)$, $C(7;5;8)$, $D(3;-1;2)$, $E(-4;-1;4)$. Ba điểm nào dưới đây thẳng hàng?}
  {$A, C, D$.}
  {$B, C, E$.}{$A, D, E$.}{$B, C, D$.}

\cauTN[id=C2-B8-D09-NB-009,muc=NB,loigiai={Ba điểm thẳng hàng khi hai vectơ tạo bởi chúng cùng phương. Ta có $\overrightarrow{BC}=(-3;-2;1)$ và $\overrightarrow{BE}=(-6;-4;2)$ cùng phương nên $B$, $C$, $E$ thẳng hàng. Các bộ còn lại không thoả mãn.}]{Trong không gian $Oxyz$, cho $A(-1;0;1)$, $B(5;5;2)$, $C(2;3;3)$, $D(-2;0;0)$, $E(-1;1;4)$. Ba điểm nào dưới đây thẳng hàng?}
  {$B, C, E$.}
  {$A, D, E$.}{$A, B, D$.}{$C, D, E$.}

\cauTN[id=C2-B8-D09-NB-010,muc=NB,loigiai={Ba điểm thẳng hàng khi hai vectơ tạo bởi chúng cùng phương. Ta có $\overrightarrow{AB}=(-2;-3;0)$ và $\overrightarrow{AC}=(2;3;0)$ cùng phương nên $A$, $B$, $C$ thẳng hàng. Các bộ còn lại không thoả mãn.}]{Trong không gian $Oxyz$, cho $A(-3;-4;0)$, $B(-5;-7;0)$, $C(-1;-1;0)$, $D(3;4;-2)$, $E(-4;3;4)$. Ba điểm nào dưới đây thẳng hàng?}
  {$A, B, C$.}
  {$A, D, E$.}{$A, B, D$.}{$A, C, E$.}

\cauTN[id=C2-B8-D09-NB-011,muc=NB,loigiai={Ba điểm thẳng hàng khi hai vectơ tạo bởi chúng cùng phương. Ta có $\overrightarrow{AC}=(0;8;-8)$ và $\overrightarrow{AE}=(0;4;-4)$ cùng phương nên $A$, $C$, $E$ thẳng hàng. Các bộ còn lại không thoả mãn.}]{Trong không gian $Oxyz$, cho $A(-3;-1;1)$, $B(1;4;-4)$, $C(-3;7;-7)$, $D(2;3;0)$, $E(-3;3;-3)$. Ba điểm nào dưới đây thẳng hàng?}
  {$A, C, E$.}
  {$A, B, E$.}{$C, D, E$.}{$A, B, D$.}

\cauTN[id=C2-B8-D09-NB-012,muc=NB,loigiai={Ba điểm thẳng hàng khi hai vectơ tạo bởi chúng cùng phương. Ta có $\overrightarrow{BC}=(-2;-1;5)$ và $\overrightarrow{BE}=(-4;-2;10)$ cùng phương nên $B$, $C$, $E$ thẳng hàng. Các bộ còn lại không thoả mãn.}]{Trong không gian $Oxyz$, cho $A(-2;0;3)$, $B(2;3;-7)$, $C(0;2;-2)$, $D(-4;-3;0)$, $E(-2;1;3)$. Ba điểm nào dưới đây thẳng hàng?}
  {$B, C, E$.}
  {$C, D, E$.}{$A, C, D$.}{$A, B, C$.}

\cauTN[id=C2-B8-D09-NB-013,muc=NB,loigiai={Hai vectơ cùng phương khi vectơ này bằng $k$ lần vectơ kia. Ta có $(-2;4;2)=2\cdot(-1;2;1)$ nên cùng phương với $\vec{a}$; các vectơ còn lại không tỉ lệ.}]{Vectơ nào dưới đây cùng phương với vectơ $\vec{a}=(-1;2;1)$?}
  {$(-2;4;2)$.}
  {$(-2;4;4)$.}{$(-1;4;2)$.}{$(-3;4;2)$.}

\cauTN[id=C2-B8-D09-NB-014,muc=NB,loigiai={Hai vectơ cùng phương khi vectơ này bằng $k$ lần vectơ kia. Ta có $(-6;-2;4)=-2\cdot(3;1;-2)$ nên cùng phương với $\vec{a}$; các vectơ còn lại không tỉ lệ.}]{Vectơ nào dưới đây cùng phương với vectơ $\vec{a}=(3;1;-2)$?}
  {$(-6;-2;4)$.}
  {$(-6;-2;3)$.}{$(-6;-2;-4)$.}{$(-8;-2;4)$.}

\cauTN[id=C2-B8-D09-NB-015,muc=NB,loigiai={Hai vectơ cùng phương khi vectơ này bằng $k$ lần vectơ kia. Ta có $(-9;0;-6)=-3\cdot(3;0;2)$ nên cùng phương với $\vec{a}$; các vectơ còn lại không tỉ lệ.}]{Vectơ nào dưới đây cùng phương với vectơ $\vec{a}=(3;0;2)$?}
  {$(-9;0;-6)$.}
  {$(-7;0;-6)$.}{$(-9;0;-5)$.}{$(-9;0;-4)$.}

\cauTN[id=C2-B8-D09-NB-016,muc=NB,loigiai={Hai vectơ cùng phương khi vectơ này bằng $k$ lần vectơ kia. Ta có $(-3;-2;-2)=-1\cdot(3;2;2)$ nên cùng phương với $\vec{a}$; các vectơ còn lại không tỉ lệ.}]{Vectơ nào dưới đây cùng phương với vectơ $\vec{a}=(3;2;2)$?}
  {$(-3;-2;-2)$.}
  {$(-3;-4;-2)$.}{$(-2;-3;-2)$.}{$(-3;-2;-4)$.}

\cauTN[id=C2-B8-D09-NB-017,muc=NB,loigiai={Hai vectơ cùng phương khi vectơ này bằng $k$ lần vectơ kia. Ta có $(-6;-2;-2)=-2\cdot(3;1;1)$ nên cùng phương với $\vec{a}$; các vectơ còn lại không tỉ lệ.}]{Vectơ nào dưới đây cùng phương với vectơ $\vec{a}=(3;1;1)$?}
  {$(-6;-2;-2)$.}
  {$(-5;-2;-2)$.}{$(-6;-1;-2)$.}{$(-6;-4;-2)$.}

\cauTN[id=C2-B8-D09-NB-018,muc=NB,loigiai={Hai vectơ cùng phương khi vectơ này bằng $k$ lần vectơ kia. Ta có $(3;0;0)=-3\cdot(-1;0;0)$ nên cùng phương với $\vec{a}$; các vectơ còn lại không tỉ lệ.}]{Vectơ nào dưới đây cùng phương với vectơ $\vec{a}=(-1;0;0)$?}
  {$(3;0;0)$.}
  {$(3;0;1)$.}{$(3;1;0)$.}{$(0;3;0)$.}

\cauTN[id=C2-B8-D09-NB-019,muc=NB,loigiai={Hai vectơ cùng phương khi vectơ này bằng $k$ lần vectơ kia. Ta có $(0;1;2)=-1\cdot(0;-1;-2)$ nên cùng phương với $\vec{a}$; các vectơ còn lại không tỉ lệ.}]{Vectơ nào dưới đây cùng phương với vectơ $\vec{a}=(0;-1;-2)$?}
  {$(0;1;2)$.}
  {$(0;-1;2)$.}{$(1;0;2)$.}{$(0;3;2)$.}

\cauTN[id=C2-B8-D09-NB-020,muc=NB,loigiai={Hai vectơ cùng phương khi vectơ này bằng $k$ lần vectơ kia. Ta có $(0;-6;9)=-3\cdot(0;2;-3)$ nên cùng phương với $\vec{a}$; các vectơ còn lại không tỉ lệ.}]{Vectơ nào dưới đây cùng phương với vectơ $\vec{a}=(0;2;-3)$?}
  {$(0;-6;9)$.}
  {$(0;-8;9)$.}{$(1;-6;9)$.}{$(0;-4;9)$.}

\cauTN[id=C2-B8-D09-NB-021,muc=NB,loigiai={Hai vectơ cùng phương khi vectơ này bằng $k$ lần vectơ kia. Ta có $(-1;-1;2)=-1\cdot(1;1;-2)$ nên cùng phương với $\vec{a}$; các vectơ còn lại không tỉ lệ.}]{Vectơ nào dưới đây cùng phương với vectơ $\vec{a}=(1;1;-2)$?}
  {$(-1;-1;2)$.}
  {$(1;-1;2)$.}{$(-1;1;2)$.}{$(-1;2;-1)$.}

\cauTN[id=C2-B8-D09-NB-022,muc=NB,loigiai={Hai vectơ cùng phương khi vectơ này bằng $k$ lần vectơ kia. Ta có $(-6;0;-3)=-3\cdot(2;0;1)$ nên cùng phương với $\vec{a}$; các vectơ còn lại không tỉ lệ.}]{Vectơ nào dưới đây cùng phương với vectơ $\vec{a}=(2;0;1)$?}
  {$(-6;0;-3)$.}
  {$(-6;0;-1)$.}{$(-6;2;-3)$.}{$(-7;0;-3)$.}

\cauTN[id=C2-B8-D09-NB-023,muc=NB,loigiai={Hai vectơ cùng phương khi vectơ này bằng $k$ lần vectơ kia. Ta có $(-2;-6;-2)=2\cdot(-1;-3;-1)$ nên cùng phương với $\vec{a}$; các vectơ còn lại không tỉ lệ.}]{Vectơ nào dưới đây cùng phương với vectơ $\vec{a}=(-1;-3;-1)$?}
  {$(-2;-6;-2)$.}
  {$(-2;-6;0)$.}{$(-3;-6;-2)$.}{$(-2;-6;-3)$.}

\cauTN[id=C2-B8-D09-NB-024,muc=NB,loigiai={Hai vectơ cùng phương khi vectơ này bằng $k$ lần vectơ kia. Ta có $(0;-9;6)=-3\cdot(0;3;-2)$ nên cùng phương với $\vec{a}$; các vectơ còn lại không tỉ lệ.}]{Vectơ nào dưới đây cùng phương với vectơ $\vec{a}=(0;3;-2)$?}
  {$(0;-9;6)$.}
  {$(0;-9;7)$.}{$(0;-11;6)$.}{$(-1;-9;6)$.}

\cauTN[id=C2-B8-D09-NB-025,muc=NB,loigiai={Tứ giác $ABCD$ là hình bình hành khi $\overrightarrow{AB}=\overrightarrow{DC}$, suy ra $D=A+C-B=(-3;1;-5)$.}]{Trong không gian $Oxyz$, cho $A(-1;-3;1)$, $B(2;-1;3)$, $C(0;3;-3)$. Tìm tọa độ điểm $D$ sao cho $ABCD$ là hình bình hành.}
  {$(-3;1;-5)$.}
  {$(3;-1;5)$.}{$(3;5;-1)$.}{$(-3;-1;-5)$.}

\cauTN[id=C2-B8-D09-NB-026,muc=NB,loigiai={Tứ giác $ABCD$ là hình bình hành khi $\overrightarrow{AB}=\overrightarrow{DC}$, suy ra $D=A+C-B=(5;0;4)$.}]{Trong không gian $Oxyz$, cho $A(2;-2;2)$, $B(-3;-3;1)$, $C(0;-1;3)$. Tìm tọa độ điểm $D$ sao cho $ABCD$ là hình bình hành.}
  {$(5;0;4)$.}
  {$(5;-1;4)$.}{$(5;0;6)$.}{$(-5;-2;2)$.}

\cauTN[id=C2-B8-D09-NB-027,muc=NB,loigiai={Tứ giác $ABCD$ là hình bình hành khi $\overrightarrow{AB}=\overrightarrow{DC}$, suy ra $D=A+C-B=(4;3;0)$.}]{Trong không gian $Oxyz$, cho $A(3;0;1)$, $B(1;-3;2)$, $C(2;0;1)$. Tìm tọa độ điểm $D$ sao cho $ABCD$ là hình bình hành.}
  {$(4;3;0)$.}
  {$(2;3;0)$.}{$(0;3;4)$.}{$(4;0;3)$.}

\cauTN[id=C2-B8-D09-NB-028,muc=NB,loigiai={Tứ giác $ABCD$ là hình bình hành khi $\overrightarrow{AB}=\overrightarrow{DC}$, suy ra $D=A+C-B=(2;-6;-4)$.}]{Trong không gian $Oxyz$, cho $A(2;-3;0)$, $B(-1;1;2)$, $C(-1;-2;-2)$. Tìm tọa độ điểm $D$ sao cho $ABCD$ là hình bình hành.}
  {$(2;-6;-4)$.}
  {$(2;6;-4)$.}{$(2;-6;-5)$.}{$(2;0;4)$.}

\cauTN[id=C2-B8-D09-NB-029,muc=NB,loigiai={Tứ giác $ABCD$ là hình bình hành khi $\overrightarrow{AB}=\overrightarrow{DC}$, suy ra $D=A+C-B=(-3;1;-1)$.}]{Trong không gian $Oxyz$, cho $A(-2;-1;-2)$, $B(1;-3;-3)$, $C(0;-1;-2)$. Tìm tọa độ điểm $D$ sao cho $ABCD$ là hình bình hành.}
  {$(-3;1;-1)$.}
  {$(3;-1;1)$.}{$(-3;1;1)$.}{$(-1;1;-3)$.}

\cauTN[id=C2-B8-D09-NB-030,muc=NB,loigiai={Tứ giác $ABCD$ là hình bình hành khi $\overrightarrow{AB}=\overrightarrow{DC}$, suy ra $D=A+C-B=(-1;7;1)$.}]{Trong không gian $Oxyz$, cho $A(3;3;0)$, $B(3;-2;0)$, $C(-1;2;1)$. Tìm tọa độ điểm $D$ sao cho $ABCD$ là hình bình hành.}
  {$(-1;7;1)$.}
  {$(7;-1;-1)$.}{$(-1;8;1)$.}{$(1;7;-1)$.}

\cauTN[id=C2-B8-D09-NB-031,muc=NB,loigiai={Tứ giác $ABCD$ là hình bình hành khi $\overrightarrow{AB}=\overrightarrow{DC}$, suy ra $D=A+C-B=(-2;5;-4)$.}]{Trong không gian $Oxyz$, cho $A(-3;1;2)$, $B(-2;-1;3)$, $C(-1;3;-3)$. Tìm tọa độ điểm $D$ sao cho $ABCD$ là hình bình hành.}
  {$(-2;5;-4)$.}
  {$(-4;5;-4)$.}{$(-2;7;-4)$.}{$(2;-5;4)$.}

\cauTN[id=C2-B8-D09-NB-032,muc=NB,loigiai={Tứ giác $ABCD$ là hình bình hành khi $\overrightarrow{AB}=\overrightarrow{DC}$, suy ra $D=A+C-B=(-3;-2;-2)$.}]{Trong không gian $Oxyz$, cho $A(-1;-1;0)$, $B(-1;3;2)$, $C(-3;2;0)$. Tìm tọa độ điểm $D$ sao cho $ABCD$ là hình bình hành.}
  {$(-3;-2;-2)$.}
  {$(-3;-2;2)$.}{$(-3;0;-2)$.}{$(-2;-2;-2)$.}

\cauTN[id=C2-B8-D09-NB-033,muc=NB,loigiai={Tứ giác $ABCD$ là hình bình hành khi $\overrightarrow{AB}=\overrightarrow{DC}$, suy ra $D=A+C-B=(-4;4;-2)$.}]{Trong không gian $Oxyz$, cho $A(-3;3;0)$, $B(-2;1;0)$, $C(-3;2;-2)$. Tìm tọa độ điểm $D$ sao cho $ABCD$ là hình bình hành.}
  {$(-4;4;-2)$.}
  {$(-4;4;0)$.}{$(-4;4;-4)$.}{$(4;4;-2)$.}

\cauTN[id=C2-B8-D09-NB-034,muc=NB,loigiai={Tứ giác $ABCD$ là hình bình hành khi $\overrightarrow{AB}=\overrightarrow{DC}$, suy ra $D=A+C-B=(4;-4;0)$.}]{Trong không gian $Oxyz$, cho $A(0;1;1)$, $B(-2;3;3)$, $C(2;-2;2)$. Tìm tọa độ điểm $D$ sao cho $ABCD$ là hình bình hành.}
  {$(4;-4;0)$.}
  {$(4;0;-4)$.}{$(4;-5;0)$.}{$(-4;6;2)$.}

\cauTN[id=C2-B8-D09-NB-035,muc=NB,loigiai={Tứ giác $ABCD$ là hình bình hành khi $\overrightarrow{AB}=\overrightarrow{DC}$, suy ra $D=A+C-B=(0;6;-7)$.}]{Trong không gian $Oxyz$, cho $A(2;0;-1)$, $B(-1;-3;3)$, $C(-3;3;-3)$. Tìm tọa độ điểm $D$ sao cho $ABCD$ là hình bình hành.}
  {$(0;6;-7)$.}
  {$(0;5;-7)$.}{$(0;6;-5)$.}{$(0;-6;-7)$.}

\cauTN[id=C2-B8-D09-NB-036,muc=NB,loigiai={Tứ giác $ABCD$ là hình bình hành khi $\overrightarrow{AB}=\overrightarrow{DC}$, suy ra $D=A+C-B=(-6;-1;3)$.}]{Trong không gian $Oxyz$, cho $A(-3;-2;3)$, $B(1;-3;0)$, $C(-2;-2;0)$. Tìm tọa độ điểm $D$ sao cho $ABCD$ là hình bình hành.}
  {$(-6;-1;3)$.}
  {$(-6;1;3)$.}{$(-6;-1;1)$.}{$(-6;-1;-3)$.}

\cauDS[id=C2-B8-D09-NB-037,muc=NB,loigiai={Kiểm tra từng mệnh đề bằng công thức tọa độ. a) Đúng. b) Đúng. c) Sai. d) Sai.}]{Trong không gian $Oxyz$, cho $A(2;-2;-2)$, $B(3;1;-3)$. Xét tính đúng sai của các mệnh đề sau:}{%
  \y{Vectơ $\overrightarrow{AB}$ có tọa độ $(1;3;-1)$.}{d}
  \y{Độ dài đoạn thẳng $AB$ bằng $\sqrt{11}$.}{d}
  \y{Với $A(4;-1;-1)$, $B(0;-1;-4)$, $C(12;-1;5)$ thẳng hàng thì $\overrightarrow{AC}=-\dfrac{1}{2}\,\overrightarrow{AB}$.}{s}
  \y{Vectơ $\overrightarrow{AB}$ có tọa độ $(1;3;0)$.}{s}}

\cauDS[id=C2-B8-D09-NB-038,muc=NB,loigiai={Kiểm tra từng mệnh đề bằng công thức tọa độ. a) Đúng. b) Sai. c) Đúng. d) Sai.}]{Trong không gian $Oxyz$, cho $A(-3;0;1)$, $B(2;2;2)$. Xét tính đúng sai của các mệnh đề sau:}{%
  \y{Độ dài đoạn thẳng $AB$ bằng $\sqrt{30}$.}{d}
  \y{Độ dài đoạn thẳng $AB$ bằng $\sqrt{29}$.}{s}
  \y{Vectơ $\overrightarrow{AB}$ có tọa độ $(5;2;1)$.}{d}
  \y{Với $A(3;-3;0)$, $B(1;-4;-3)$, $C(-3;-6;-9)$ thẳng hàng thì $\overrightarrow{AC}=7\,\overrightarrow{AB}$.}{s}}

\cauDS[id=C2-B8-D09-NB-039,muc=NB,loigiai={Kiểm tra từng mệnh đề bằng công thức tọa độ. a) Đúng. b) Sai. c) Đúng. d) Sai.}]{Trong không gian $Oxyz$, cho $A(1;3;1)$, $B(0;1;1)$, $C(-3;-1;0)$. Xét tính đúng sai của các mệnh đề sau:}{%
  \y{Độ dài đoạn thẳng $AB$ bằng $\sqrt{5}$.}{d}
  \y{Độ dài đoạn thẳng $AB$ bằng $\sqrt{6}$.}{s}
  \y{Điểm $D$ để $ABCD$ là hình bình hành có tọa độ $(-2;1;0)$.}{d}
  \y{Vectơ $\overrightarrow{AB}$ có tọa độ $(-3;-2;0)$.}{s}}

\cauDS[id=C2-B8-D09-NB-040,muc=NB,loigiai={Kiểm tra từng mệnh đề bằng công thức tọa độ. a) Đúng. b) Sai. c) Sai. d) Đúng.}]{Trong không gian $Oxyz$, cho $A(-2;1;-3)$, $B(-3;-2;-1)$, $C(-2;3;-2)$. Xét tính đúng sai của các mệnh đề sau:}{%
  \y{Với $A(4;-3;-4)$, $B(1;-3;-3)$, $C(-2;-3;-2)$ thẳng hàng thì $\overrightarrow{AC}=2\,\overrightarrow{AB}$.}{d}
  \y{Vectơ $\overrightarrow{AB}$ có tọa độ $(-1;-1;2)$.}{s}
  \y{Điểm $D$ để $ABCD$ là hình bình hành có tọa độ $(1;6;-4)$.}{s}
  \y{Với $A(-1;-3;-4)$, $B(3;-4;-1)$, $C(-9;-1;-10)$ thẳng hàng thì $\overrightarrow{AC}=-2\,\overrightarrow{AB}$.}{d}}

\cauDS[id=C2-B8-D09-NB-041,muc=NB,loigiai={Kiểm tra từng mệnh đề bằng công thức tọa độ. a) Đúng. b) Sai. c) Đúng. d) Sai.}]{Trong không gian $Oxyz$, cho $A(-2;3;-2)$, $B(1;2;-1)$, $C(3;0;3)$. Xét tính đúng sai của các mệnh đề sau:}{%
  \y{Với $A(3;-3;-1)$, $B(-4;3;2)$, $C(-18;15;8)$ thẳng hàng thì $\overrightarrow{AC}=3\,\overrightarrow{AB}$.}{d}
  \y{Điểm $D$ để $ABCD$ là hình bình hành có tọa độ $(0;1;4)$.}{s}
  \y{Độ dài đoạn thẳng $AB$ bằng $\sqrt{11}$.}{d}
  \y{Vectơ $\overrightarrow{AB}$ có tọa độ $(3;-1;2)$.}{s}}

\cauDS[id=C2-B8-D09-NB-042,muc=NB,loigiai={Kiểm tra từng mệnh đề bằng công thức tọa độ. a) Đúng. b) Sai. c) Sai. d) Sai.}]{Trong không gian $Oxyz$, cho $A(-3;-3;3)$, $B(0;-3;-3)$, $C(-3;-3;2)$. Xét tính đúng sai của các mệnh đề sau:}{%
  \y{Với $A(-4;-3;1)$, $B(-4;-2;4)$, $C(-4;-5;-5)$ thẳng hàng thì $\overrightarrow{AC}=-2\,\overrightarrow{AB}$.}{d}
  \y{Điểm $D$ để $ABCD$ là hình bình hành có tọa độ $(-6;3;8)$.}{s}
  \y{Độ dài đoạn thẳng $AB$ bằng $\sqrt{46}$.}{s}
  \y{Với $A(4;3;-3)$, $B(0;-3;-1)$, $C(12;15;-7)$ thẳng hàng thì $\overrightarrow{AC}=-1\,\overrightarrow{AB}$.}{s}}

\cauDS[id=C2-B8-D09-NB-043,muc=NB,loigiai={Kiểm tra từng mệnh đề bằng công thức tọa độ. a) Đúng. b) Sai. c) Đúng. d) Đúng.}]{Trong không gian $Oxyz$, cho $A(-3;2;3)$, $B(1;-2;0)$, $C(-2;3;-1)$. Xét tính đúng sai của các mệnh đề sau:}{%
  \y{Độ dài đoạn thẳng $AB$ bằng $\sqrt{41}$.}{d}
  \y{Điểm $D$ để $ABCD$ là hình bình hành có tọa độ $(0;-3;4)$.}{s}
  \y{Vectơ $\overrightarrow{AB}$ có tọa độ $(4;-4;-3)$.}{d}
  \y{Với $A(-2;-1;2)$, $B(-1;2;1)$, $C(0;5;0)$ thẳng hàng thì $\overrightarrow{AC}=2\,\overrightarrow{AB}$.}{d}}

\cauDS[id=C2-B8-D09-NB-044,muc=NB,loigiai={Kiểm tra từng mệnh đề bằng công thức tọa độ. a) Sai. b) Đúng. c) Sai. d) Đúng.}]{Trong không gian $Oxyz$, cho $A(3;2;-3)$, $B(1;3;3)$, $C(2;0;-3)$. Xét tính đúng sai của các mệnh đề sau:}{%
  \y{Vectơ $\overrightarrow{AB}$ có tọa độ $(4;5;0)$.}{s}
  \y{Vectơ $\overrightarrow{AB}$ có tọa độ $(-2;1;6)$.}{d}
  \y{Điểm $D$ để $ABCD$ là hình bình hành có tọa độ $(6;-1;-9)$.}{s}
  \y{Độ dài đoạn thẳng $AB$ bằng $\sqrt{41}$.}{d}}

\cauDS[id=C2-B8-D09-NB-045,muc=NB,loigiai={Kiểm tra từng mệnh đề bằng công thức tọa độ. a) Sai. b) Sai. c) Đúng. d) Đúng.}]{Trong không gian $Oxyz$, cho $A(1;0;-3)$, $B(-3;-1;3)$, $C(0;3;3)$. Xét tính đúng sai của các mệnh đề sau:}{%
  \y{Điểm $D$ để $ABCD$ là hình bình hành có tọa độ $(-4;4;-3)$.}{s}
  \y{Độ dài đoạn thẳng $AB$ bằng $3\sqrt{6}$.}{s}
  \y{Vectơ $\overrightarrow{AB}$ có tọa độ $(-4;-1;6)$.}{d}
  \y{Độ dài đoạn thẳng $AB$ bằng $\sqrt{53}$.}{d}}

\cauDS[id=C2-B8-D09-NB-046,muc=NB,loigiai={Kiểm tra từng mệnh đề bằng công thức tọa độ. a) Đúng. b) Sai. c) Sai. d) Đúng.}]{Trong không gian $Oxyz$, cho $A(-2;-3;1)$, $B(-2;2;1)$. Xét tính đúng sai của các mệnh đề sau:}{%
  \y{Với $A(4;-1;-1)$, $B(3;1;3)$, $C(6;-5;-9)$ thẳng hàng thì $\overrightarrow{AC}=-2\,\overrightarrow{AB}$.}{d}
  \y{Độ dài đoạn thẳng $AB$ bằng $2\sqrt{6}$.}{s}
  \y{Vectơ $\overrightarrow{AB}$ có tọa độ $(0;5;-2)$.}{s}
  \y{Với $A(-1;4;2)$, $B(1;-4;1)$, $C(3;-12;0)$ thẳng hàng thì $\overrightarrow{AC}=2\,\overrightarrow{AB}$.}{d}}

\cauDS[id=C2-B8-D09-NB-047,muc=NB,loigiai={Kiểm tra từng mệnh đề bằng công thức tọa độ. a) Sai. b) Đúng. c) Đúng. d) Đúng.}]{Trong không gian $Oxyz$, cho $A(-2;1;0)$, $B(1;-3;2)$. Xét tính đúng sai của các mệnh đề sau:}{%
  \y{Với $A(-2;-3;0)$, $B(-4;2;4)$, $C(2;-13;-8)$ thẳng hàng thì $\overrightarrow{AC}=-5\,\overrightarrow{AB}$.}{s}
  \y{Với $A(-3;1;-1)$, $B(3;-4;3)$, $C(9;-9;7)$ thẳng hàng thì $\overrightarrow{AC}=2\,\overrightarrow{AB}$.}{d}
  \y{Độ dài đoạn thẳng $AB$ bằng $\sqrt{29}$.}{d}
  \y{Vectơ $\overrightarrow{AB}$ có tọa độ $(3;-4;2)$.}{d}}

\cauDS[id=C2-B8-D09-NB-048,muc=NB,loigiai={Kiểm tra từng mệnh đề bằng công thức tọa độ. a) Sai. b) Đúng. c) Đúng. d) Sai.}]{Trong không gian $Oxyz$, cho $A(-2;0;0)$, $B(-2;3;3)$, $C(3;-2;2)$. Xét tính đúng sai của các mệnh đề sau:}{%
  \y{Điểm $D$ để $ABCD$ là hình bình hành có tọa độ $(2;-5;-1)$.}{s}
  \y{Với $A(-4;-3;-4)$, $B(-2;2;4)$, $C(-6;-8;-12)$ thẳng hàng thì $\overrightarrow{AC}=-1\,\overrightarrow{AB}$.}{d}
  \y{Vectơ $\overrightarrow{AB}$ có tọa độ $(0;3;3)$.}{d}
  \y{Độ dài đoạn thẳng $AB$ bằng $\sqrt{22}$.}{s}}

\cauTLN[id=C2-B8-D09-NB-049,muc=NB,loigiai={Tứ giác $ABCD$ là hình bình hành khi $\overrightarrow{AB}=\overrightarrow{DC}$, suy ra $D=A+C-B=(2;4;4)$.}]{Trong không gian $Oxyz$, cho $A(2;3;1)$, $B(-3;0;-1)$, $C(-3;1;2)$. Tìm điểm $D$ sao cho $ABCD$ là hình bình hành. Tính $x_D+y_D+z_D$.}{10}

\cauTLN[id=C2-B8-D09-NB-050,muc=NB,loigiai={Tứ giác $ABCD$ là hình bình hành khi $\overrightarrow{AB}=\overrightarrow{DC}$, suy ra $D=A+C-B=(-4;0;3)$.}]{Trong không gian $Oxyz$, cho $A(-2;0;2)$, $B(0;1;0)$, $C(-2;1;1)$. Tìm điểm $D$ sao cho $ABCD$ là hình bình hành. Tính $x_D+y_D+z_D$.}{-1}

\cauTLN[id=C2-B8-D09-NB-051,muc=NB,loigiai={Tứ giác $ABCD$ là hình bình hành khi $\overrightarrow{AB}=\overrightarrow{DC}$, suy ra $D=A+C-B=(4;6;-1)$.}]{Trong không gian $Oxyz$, cho $A(0;1;-3)$, $B(-2;-3;-2)$, $C(2;2;0)$. Tìm điểm $D$ sao cho $ABCD$ là hình bình hành. Tính $x_D+y_D+z_D$.}{9}

\cauTLN[id=C2-B8-D09-NB-052,muc=NB,loigiai={Tứ giác $ABCD$ là hình bình hành khi $\overrightarrow{AB}=\overrightarrow{DC}$, suy ra $D=A+C-B=(3;0;1)$.}]{Trong không gian $Oxyz$, cho $A(1;0;-1)$, $B(-1;-1;-3)$, $C(1;-1;-1)$. Tìm điểm $D$ sao cho $ABCD$ là hình bình hành. Tính $x_D+y_D+z_D$.}{4}

\cauTLN[id=C2-B8-D09-NB-053,muc=NB,loigiai={Ba điểm $A,B,C$ thẳng hàng nên $\overrightarrow{AC}=k\overrightarrow{AB}$. Từ hoành độ: $k=-1$, suy ra $z_C=2-1\cdot(-4)=6$.}]{Cho ba điểm $A(-4;-2;2)$, $B(1;4;-2)$, $C(-9;-8;c)$ thẳng hàng. Tính $c$.}{6}

\cauTLN[id=C2-B8-D09-NB-054,muc=NB,loigiai={Ba điểm $A,B,C$ thẳng hàng nên $\overrightarrow{AC}=k\overrightarrow{AB}$. Từ hoành độ: $k=2$, suy ra $z_C=-3+2\cdot(-1)=-5$.}]{Cho ba điểm $A(1;-4;-3)$, $B(-2;-1;-4)$, $C(-5;2;c)$ thẳng hàng. Tính $c$.}{-5}

\cauTLN[id=C2-B8-D09-NB-055,muc=NB,loigiai={Ba điểm $A,B,C$ thẳng hàng nên $\overrightarrow{AC}=k\overrightarrow{AB}$. Từ hoành độ: $k=3$, suy ra $z_C=2+3\cdot(0)=2$.}]{Cho ba điểm $A(1;-4;2)$, $B(3;-2;2)$, $C(7;2;c)$ thẳng hàng. Tính $c$.}{2}

\cauTLN[id=C2-B8-D09-NB-056,muc=NB,loigiai={Ba điểm $A,B,C$ thẳng hàng nên $\overrightarrow{AC}=k\overrightarrow{AB}$. Từ hoành độ: $k=-2$, suy ra $z_C=4-2\cdot(-8)=20$.}]{Cho ba điểm $A(3;3;4)$, $B(-3;1;-4)$, $C(15;7;c)$ thẳng hàng. Tính $c$.}{20}

\cauTLN[id=C2-B8-D09-NB-057,muc=NB,loigiai={Ta có $\overrightarrow{AB}=(-1;0;-2)$, $\overrightarrow{AC}=(-2;0;-4)$. Vì $\overrightarrow{AC}=k\overrightarrow{AB}$ nên $k=2$.}]{Cho ba điểm thẳng hàng $A(-2;-3;1)$, $B(-3;-3;-1)$, $C(-4;-3;-3)$. Biết $\overrightarrow{AC}=k\overrightarrow{AB}$. Tính $k$.}{2}

\cauTLN[id=C2-B8-D09-NB-058,muc=NB,loigiai={Ta có $\overrightarrow{AB}=(7;-1;-7)$, $\overrightarrow{AC}=(21;-3;-21)$. Vì $\overrightarrow{AC}=k\overrightarrow{AB}$ nên $k=3$.}]{Cho ba điểm thẳng hàng $A(-3;-3;3)$, $B(4;-4;-4)$, $C(18;-6;-18)$. Biết $\overrightarrow{AC}=k\overrightarrow{AB}$. Tính $k$.}{3}

\cauTLN[id=C2-B8-D09-NB-059,muc=NB,loigiai={Ta có $\overrightarrow{AB}=(-1;-1;-3)$, $\overrightarrow{AC}=(-3;-3;-9)$. Vì $\overrightarrow{AC}=k\overrightarrow{AB}$ nên $k=3$.}]{Cho ba điểm thẳng hàng $A(-2;-1;1)$, $B(-3;-2;-2)$, $C(-5;-4;-8)$. Biết $\overrightarrow{AC}=k\overrightarrow{AB}$. Tính $k$.}{3}

\cauTLN[id=C2-B8-D09-NB-060,muc=NB,loigiai={Ta có $\overrightarrow{AB}=(-2;-2;2)$, $\overrightarrow{AC}=(-4;-4;4)$. Vì $\overrightarrow{AC}=k\overrightarrow{AB}$ nên $k=2$.}]{Cho ba điểm thẳng hàng $A(0;2;1)$, $B(-2;0;3)$, $C(-4;-2;5)$. Biết $\overrightarrow{AC}=k\overrightarrow{AB}$. Tính $k$.}{2}

\cauTN[id=C2-B8-D09-TH-001,muc=TH,loigiai={Ba điểm $A,B,C$ thẳng hàng nên $\overrightarrow{AC}=k\overrightarrow{AB}$. Từ hoành độ: $k=-2$, suy ra $z_C=-3-2\cdot(6)=-15$.}]{Cho ba điểm $A(1;2;-3)$, $B(-2;1;3)$, $C(7;4;c)$ thẳng hàng. Giá trị của $c$ bằng}
  {$-15$.}
  {$-16$.}{$15$.}{$-13$.}

\cauTN[id=C2-B8-D09-TH-002,muc=TH,loigiai={Ba điểm $A,B,C$ thẳng hàng nên $\overrightarrow{AC}=k\overrightarrow{AB}$. Từ hoành độ: $k=2$, suy ra $z_C=0+2\cdot(-4)=-8$.}]{Cho ba điểm $A(-3;-1;0)$, $B(-1;3;-4)$, $C(1;7;c)$ thẳng hàng. Giá trị của $c$ bằng}
  {$-8$.}
  {$-24$.}{$-15$.}{$-9$.}

\cauTN[id=C2-B8-D09-TH-003,muc=TH,loigiai={Ba điểm $A,B,C$ thẳng hàng nên $\overrightarrow{AC}=k\overrightarrow{AB}$. Từ hoành độ: $k=-1$, suy ra $z_C=1-1\cdot(-5)=6$.}]{Cho ba điểm $A(-2;1;1)$, $B(-3;-3;-4)$, $C(-1;5;c)$ thẳng hàng. Giá trị của $c$ bằng}
  {$6$.}
  {$18$.}{$7$.}{$3$.}

\cauTN[id=C2-B8-D09-TH-004,muc=TH,loigiai={Ba điểm $A,B,C$ thẳng hàng nên $\overrightarrow{AC}=k\overrightarrow{AB}$. Từ hoành độ: $k=3$, suy ra $z_C=-4+3\cdot(4)=8$.}]{Cho ba điểm $A(4;-3;-4)$, $B(1;0;0)$, $C(-5;6;c)$ thẳng hàng. Giá trị của $c$ bằng}
  {$8$.}
  {$11$.}{$-8$.}{$16$.}

\cauTN[id=C2-B8-D09-TH-005,muc=TH,loigiai={Ba điểm $A,B,C$ thẳng hàng nên $\overrightarrow{AC}=k\overrightarrow{AB}$. Từ hoành độ: $k=-2$, suy ra $z_C=1-2\cdot(-5)=11$.}]{Cho ba điểm $A(-2;0;1)$, $B(3;-4;-4)$, $C(-12;8;c)$ thẳng hàng. Giá trị của $c$ bằng}
  {$11$.}
  {$9$.}{$8$.}{$13$.}

\cauTN[id=C2-B8-D09-TH-006,muc=TH,loigiai={Ba điểm $A,B,C$ thẳng hàng nên $\overrightarrow{AC}=k\overrightarrow{AB}$. Từ hoành độ: $k=-1$, suy ra $z_C=3-1\cdot(-4)=7$.}]{Cho ba điểm $A(-4;2;3)$, $B(-1;-1;-1)$, $C(-7;5;c)$ thẳng hàng. Giá trị của $c$ bằng}
  {$7$.}
  {$-7$.}{$5$.}{$14$.}

\cauTN[id=C2-B8-D09-TH-007,muc=TH,loigiai={Ba điểm $A,B,C$ thẳng hàng nên $\overrightarrow{AC}=k\overrightarrow{AB}$. Từ hoành độ: $k=2$, suy ra $z_C=0+2\cdot(3)=6$.}]{Cho ba điểm $A(-1;-3;0)$, $B(3;-3;3)$, $C(7;-3;c)$ thẳng hàng. Giá trị của $c$ bằng}
  {$6$.}
  {$-6$.}{$3$.}{$18$.}

\cauTN[id=C2-B8-D09-TH-008,muc=TH,loigiai={Ba điểm $A,B,C$ thẳng hàng nên $\overrightarrow{AC}=k\overrightarrow{AB}$. Từ hoành độ: $k=3$, suy ra $z_C=-3+3\cdot(5)=12$.}]{Cho ba điểm $A(4;-3;-3)$, $B(-1;-3;2)$, $C(-11;-3;c)$ thẳng hàng. Giá trị của $c$ bằng}
  {$12$.}
  {$9$.}{$10$.}{$15$.}

\cauTN[id=C2-B8-D09-TH-009,muc=TH,loigiai={Ba điểm $A,B,C$ thẳng hàng nên $\overrightarrow{AC}=k\overrightarrow{AB}$. Từ hoành độ: $k=-1$, suy ra $z_C=-4-1\cdot(5)=-9$.}]{Cho ba điểm $A(0;1;-4)$, $B(3;4;1)$, $C(-3;-2;c)$ thẳng hàng. Giá trị của $c$ bằng}
  {$-9$.}
  {$-8$.}{$-6$.}{$9$.}

\cauTN[id=C2-B8-D09-TH-010,muc=TH,loigiai={Ba điểm $A,B,C$ thẳng hàng nên $\overrightarrow{AC}=k\overrightarrow{AB}$. Từ hoành độ: $k=3$, suy ra $z_C=2+3\cdot(-6)=-16$.}]{Cho ba điểm $A(3;4;2)$, $B(0;-3;-4)$, $C(-6;-17;c)$ thẳng hàng. Giá trị của $c$ bằng}
  {$-16$.}
  {$-8$.}{$-13$.}{$-48$.}

\cauTN[id=C2-B8-D09-TH-011,muc=TH,loigiai={Ba điểm $A,B,C$ thẳng hàng nên $\overrightarrow{AC}=k\overrightarrow{AB}$. Từ hoành độ: $k=3$, suy ra $z_C=-2+3\cdot(4)=10$.}]{Cho ba điểm $A(-4;4;-2)$, $B(-2;1;2)$, $C(2;-5;c)$ thẳng hàng. Giá trị của $c$ bằng}
  {$10$.}
  {$30$.}{$11$.}{$20$.}

\cauTN[id=C2-B8-D09-TH-012,muc=TH,loigiai={Ba điểm $A,B,C$ thẳng hàng nên $\overrightarrow{AC}=k\overrightarrow{AB}$. Từ hoành độ: $k=2$, suy ra $z_C=-2+2\cdot(0)=-2$.}]{Cho ba điểm $A(-2;0;-2)$, $B(-1;-1;-2)$, $C(0;-2;c)$ thẳng hàng. Giá trị của $c$ bằng}
  {$-2$.}
  {$0$.}{$-3$.}{$-5$.}

\cauTN[id=C2-B8-D09-TH-013,muc=TH,loigiai={Ta có $\overrightarrow{AB}=(-3;5;4)$, $\overrightarrow{AC}=(-9;15;12)$. Vì $\overrightarrow{AC}=k\overrightarrow{AB}$ nên $k=3$.}]{Cho ba điểm thẳng hàng $A(-1;-4;-1)$, $B(-4;1;3)$, $C(-10;11;11)$. Biết $\overrightarrow{AC}=k\overrightarrow{AB}$. Giá trị của $k$ bằng}
  {$3$.}
  {$2$.}{$-3$.}{$1$.}

\cauTN[id=C2-B8-D09-TH-014,muc=TH,loigiai={Ta có $\overrightarrow{AB}=(0;-1;-5)$, $\overrightarrow{AC}=(0;1;5)$. Vì $\overrightarrow{AC}=k\overrightarrow{AB}$ nên $k=-1$.}]{Cho ba điểm thẳng hàng $A(-1;4;2)$, $B(-1;3;-3)$, $C(-1;5;7)$. Biết $\overrightarrow{AC}=k\overrightarrow{AB}$. Giá trị của $k$ bằng}
  {$-1$.}
  {$-2$.}{$2$.}{$-3$.}

\cauTN[id=C2-B8-D09-TH-015,muc=TH,loigiai={Ta có $\overrightarrow{AB}=(-1;1;6)$, $\overrightarrow{AC}=(-3;3;18)$. Vì $\overrightarrow{AC}=k\overrightarrow{AB}$ nên $k=3$.}]{Cho ba điểm thẳng hàng $A(-1;-3;-2)$, $B(-2;-2;4)$, $C(-4;0;16)$. Biết $\overrightarrow{AC}=k\overrightarrow{AB}$. Giá trị của $k$ bằng}
  {$3$.}
  {$7$.}{$6$.}{$1$.}

\cauTN[id=C2-B8-D09-TH-016,muc=TH,loigiai={Ta có $\overrightarrow{AB}=(3;-2;1)$, $\overrightarrow{AC}=(6;-4;2)$. Vì $\overrightarrow{AC}=k\overrightarrow{AB}$ nên $k=2$.}]{Cho ba điểm thẳng hàng $A(-1;-1;-1)$, $B(2;-3;0)$, $C(5;-5;1)$. Biết $\overrightarrow{AC}=k\overrightarrow{AB}$. Giá trị của $k$ bằng}
  {$2$.}
  {$-1$.}{$-2$.}{$1$.}

\cauTN[id=C2-B8-D09-TH-017,muc=TH,loigiai={Ta có $\overrightarrow{AB}=(2;1;-3)$, $\overrightarrow{AC}=(-4;-2;6)$. Vì $\overrightarrow{AC}=k\overrightarrow{AB}$ nên $k=-2$.}]{Cho ba điểm thẳng hàng $A(-3;-3;0)$, $B(-1;-2;-3)$, $C(-7;-5;6)$. Biết $\overrightarrow{AC}=k\overrightarrow{AB}$. Giá trị của $k$ bằng}
  {$-2$.}
  {$0$.}{$-3$.}{$-1$.}

\cauTN[id=C2-B8-D09-TH-018,muc=TH,loigiai={Ta có $\overrightarrow{AB}=(-1;4;1)$, $\overrightarrow{AC}=(-2;8;2)$. Vì $\overrightarrow{AC}=k\overrightarrow{AB}$ nên $k=2$.}]{Cho ba điểm thẳng hàng $A(-2;-1;-2)$, $B(-3;3;-1)$, $C(-4;7;0)$. Biết $\overrightarrow{AC}=k\overrightarrow{AB}$. Giá trị của $k$ bằng}
  {$2$.}
  {$6$.}{$5$.}{$4$.}

\cauTN[id=C2-B8-D09-TH-019,muc=TH,loigiai={Ta có $\overrightarrow{AB}=(0;-3;-1)$, $\overrightarrow{AC}=(0;-9;-3)$. Vì $\overrightarrow{AC}=k\overrightarrow{AB}$ nên $k=3$.}]{Cho ba điểm thẳng hàng $A(-1;4;3)$, $B(-1;1;2)$, $C(-1;-5;0)$. Biết $\overrightarrow{AC}=k\overrightarrow{AB}$. Giá trị của $k$ bằng}
  {$3$.}
  {$1$.}{$\dfrac{3}{2}$.}{$4$.}

\cauTN[id=C2-B8-D09-TH-020,muc=TH,loigiai={Ta có $\overrightarrow{AB}=(1;-1;0)$, $\overrightarrow{AC}=(3;-3;0)$. Vì $\overrightarrow{AC}=k\overrightarrow{AB}$ nên $k=3$.}]{Cho ba điểm thẳng hàng $A(-2;4;3)$, $B(-1;3;3)$, $C(1;1;3)$. Biết $\overrightarrow{AC}=k\overrightarrow{AB}$. Giá trị của $k$ bằng}
  {$3$.}
  {$6$.}{$5$.}{$0$.}

\cauTN[id=C2-B8-D09-TH-021,muc=TH,loigiai={Ta có $\overrightarrow{AB}=(-3;-3;0)$, $\overrightarrow{AC}=(-9;-9;0)$. Vì $\overrightarrow{AC}=k\overrightarrow{AB}$ nên $k=3$.}]{Cho ba điểm thẳng hàng $A(3;0;-3)$, $B(0;-3;-3)$, $C(-6;-9;-3)$. Biết $\overrightarrow{AC}=k\overrightarrow{AB}$. Giá trị của $k$ bằng}
  {$3$.}
  {$-3$.}{$2$.}{$0$.}

\cauTN[id=C2-B8-D09-TH-022,muc=TH,loigiai={Ta có $\overrightarrow{AB}=(7;-2;-4)$, $\overrightarrow{AC}=(14;-4;-8)$. Vì $\overrightarrow{AC}=k\overrightarrow{AB}$ nên $k=2$.}]{Cho ba điểm thẳng hàng $A(-3;1;2)$, $B(4;-1;-2)$, $C(11;-3;-6)$. Biết $\overrightarrow{AC}=k\overrightarrow{AB}$. Giá trị của $k$ bằng}
  {$2$.}
  {$1$.}{$\dfrac{1}{2}$.}{$5$.}

\cauTN[id=C2-B8-D09-TH-023,muc=TH,loigiai={Ta có $\overrightarrow{AB}=(-3;3;-2)$, $\overrightarrow{AC}=(-9;9;-6)$. Vì $\overrightarrow{AC}=k\overrightarrow{AB}$ nên $k=3$.}]{Cho ba điểm thẳng hàng $A(3;-4;-2)$, $B(0;-1;-4)$, $C(-6;5;-8)$. Biết $\overrightarrow{AC}=k\overrightarrow{AB}$. Giá trị của $k$ bằng}
  {$3$.}
  {$4$.}{$-3$.}{$9$.}

\cauTN[id=C2-B8-D09-TH-024,muc=TH,loigiai={Ta có $\overrightarrow{AB}=(3;-2;4)$, $\overrightarrow{AC}=(9;-6;12)$. Vì $\overrightarrow{AC}=k\overrightarrow{AB}$ nên $k=3$.}]{Cho ba điểm thẳng hàng $A(-2;3;-4)$, $B(1;1;0)$, $C(7;-3;8)$. Biết $\overrightarrow{AC}=k\overrightarrow{AB}$. Giá trị của $k$ bằng}
  {$3$.}
  {$-2$.}{$1$.}{$5$.}

\cauTN[id=C2-B8-D09-TH-025,muc=TH,loigiai={Hình bình hành $ABCD$ có $D=A+C-B=(-7;-7;3)$. Do đó $\overrightarrow{BD}=(-11;-9;1)$ và $BD=\sqrt{203}$.}]{Cho hình bình hành $ABCD$ với $A(1;-4;1)$, $B(4;2;2)$, $C(-4;-1;4)$. Độ dài đường chéo $BD$ bằng}
  {$\sqrt{203}$.}
  {$2\sqrt{51}$.}{$\sqrt{202}$.}{$3\sqrt{23}$.}

\cauTN[id=C2-B8-D09-TH-026,muc=TH,loigiai={Hình bình hành $ABCD$ có $D=A+C-B=(2;-4;11)$. Do đó $\overrightarrow{BD}=(7;1;16)$ và $BD=3\sqrt{34}$.}]{Cho hình bình hành $ABCD$ với $A(-2;-4;3)$, $B(-5;-5;-5)$, $C(-1;-5;3)$. Độ dài đường chéo $BD$ bằng}
  {$3\sqrt{34}$.}
  {$\sqrt{307}$.}{$\sqrt{305}$.}{$\sqrt{310}$.}

\cauTN[id=C2-B8-D09-TH-027,muc=TH,loigiai={Hình bình hành $ABCD$ có $D=A+C-B=(1;1;3)$. Do đó $\overrightarrow{BD}=(0;1;7)$ và $BD=5\sqrt{2}$.}]{Cho hình bình hành $ABCD$ với $A(-1;5;-4)$, $B(1;0;-4)$, $C(3;-4;3)$. Độ dài đường chéo $BD$ bằng}
  {$5\sqrt{2}$.}
  {$\sqrt{51}$.}{$7$.}{$3\sqrt{6}$.}

\cauTN[id=C2-B8-D09-TH-028,muc=TH,loigiai={Hình bình hành $ABCD$ có $D=A+C-B=(5;-5;-1)$. Do đó $\overrightarrow{BD}=(7;-7;-2)$ và $BD=\sqrt{102}$.}]{Cho hình bình hành $ABCD$ với $A(3;-1;-2)$, $B(-2;2;1)$, $C(0;-2;2)$. Độ dài đường chéo $BD$ bằng}
  {$\sqrt{102}$.}
  {$\sqrt{103}$.}{$\sqrt{101}$.}{$\sqrt{106}$.}

\cauTN[id=C2-B8-D09-TH-029,muc=TH,loigiai={Hình bình hành $ABCD$ có $D=A+C-B=(5;0;-7)$. Do đó $\overrightarrow{BD}=(1;-2;-11)$ và $BD=3\sqrt{14}$.}]{Cho hình bình hành $ABCD$ với $A(4;0;1)$, $B(4;2;4)$, $C(5;2;-4)$. Độ dài đường chéo $BD$ bằng}
  {$3\sqrt{14}$.}
  {$\sqrt{127}$.}{$5\sqrt{5}$.}{$\sqrt{130}$.}

\cauTN[id=C2-B8-D09-TH-030,muc=TH,loigiai={Hình bình hành $ABCD$ có $D=A+C-B=(6;-3;4)$. Do đó $\overrightarrow{BD}=(9;-3;8)$ và $BD=\sqrt{154}$.}]{Cho hình bình hành $ABCD$ với $A(5;-3;1)$, $B(-3;0;-4)$, $C(-2;0;-1)$. Độ dài đường chéo $BD$ bằng}
  {$\sqrt{154}$.}
  {$\sqrt{155}$.}{$3\sqrt{17}$.}{$\sqrt{158}$.}

\cauTN[id=C2-B8-D09-TH-031,muc=TH,loigiai={Hình bình hành $ABCD$ có $D=A+C-B=(-10;-1;-4)$. Do đó $\overrightarrow{BD}=(-12;-1;-6)$ và $BD=\sqrt{181}$.}]{Cho hình bình hành $ABCD$ với $A(-3;2;2)$, $B(2;0;2)$, $C(-5;-3;-4)$. Độ dài đường chéo $BD$ bằng}
  {$\sqrt{181}$.}
  {$\sqrt{182}$.}{$6\sqrt{5}$.}{$\sqrt{185}$.}

\cauTN[id=C2-B8-D09-TH-032,muc=TH,loigiai={Hình bình hành $ABCD$ có $D=A+C-B=(3;-8;2)$. Do đó $\overrightarrow{BD}=(5;-10;-2)$ và $BD=\sqrt{129}$.}]{Cho hình bình hành $ABCD$ với $A(-1;-2;1)$, $B(-2;2;4)$, $C(2;-4;5)$. Độ dài đường chéo $BD$ bằng}
  {$\sqrt{129}$.}
  {$\sqrt{130}$.}{$8\sqrt{2}$.}{$\sqrt{133}$.}

\cauTN[id=C2-B8-D09-TH-033,muc=TH,loigiai={Hình bình hành $ABCD$ có $D=A+C-B=(0;5;-4)$. Do đó $\overrightarrow{BD}=(3;10;-6)$ và $BD=\sqrt{145}$.}]{Cho hình bình hành $ABCD$ với $A(-4;1;-2)$, $B(-3;-5;2)$, $C(1;-1;0)$. Độ dài đường chéo $BD$ bằng}
  {$\sqrt{145}$.}
  {$\sqrt{146}$.}{$12$.}{$\sqrt{149}$.}

\cauTN[id=C2-B8-D09-TH-034,muc=TH,loigiai={Hình bình hành $ABCD$ có $D=A+C-B=(-5;3;-3)$. Do đó $\overrightarrow{BD}=(-5;7;-4)$ và $BD=3\sqrt{10}$.}]{Cho hình bình hành $ABCD$ với $A(0;2;-2)$, $B(0;-4;1)$, $C(-5;-3;0)$. Độ dài đường chéo $BD$ bằng}
  {$3\sqrt{10}$.}
  {$\sqrt{91}$.}{$\sqrt{89}$.}{$\sqrt{94}$.}

\cauTN[id=C2-B8-D09-TH-035,muc=TH,loigiai={Hình bình hành $ABCD$ có $D=A+C-B=(8;-10;3)$. Do đó $\overrightarrow{BD}=(6;-14;0)$ và $BD=2\sqrt{58}$.}]{Cho hình bình hành $ABCD$ với $A(5;-5;1)$, $B(2;4;3)$, $C(5;-1;5)$. Độ dài đường chéo $BD$ bằng}
  {$2\sqrt{58}$.}
  {$\sqrt{233}$.}{$\sqrt{231}$.}{$2\sqrt{59}$.}

\cauTN[id=C2-B8-D09-TH-036,muc=TH,loigiai={Hình bình hành $ABCD$ có $D=A+C-B=(0;13;1)$. Do đó $\overrightarrow{BD}=(-2;18;1)$ và $BD=\sqrt{329}$.}]{Cho hình bình hành $ABCD$ với $A(4;5;0)$, $B(2;-5;0)$, $C(-2;3;1)$. Độ dài đường chéo $BD$ bằng}
  {$\sqrt{329}$.}
  {$\sqrt{330}$.}{$2\sqrt{82}$.}{$3\sqrt{37}$.}

\cauTN[id=C2-B8-D09-TH-037,muc=TH,loigiai={Từ giả thiết $\overrightarrow{AM}=2\overrightarrow{AB}$ nên $M=A+2(B-A)=(1;0;6)$.}]{Trong không gian $Oxyz$, cho $A(1;2;4)$, $B(1;1;5)$. Tọa độ điểm $M$ thoả mãn $\overrightarrow{AM}=2\overrightarrow{AB}$ là}
  {$(1;0;6)$.}
  {$(1;-2;6)$.}{$(1;0;8)$.}{$(0;0;6)$.}

\cauTN[id=C2-B8-D09-TH-038,muc=TH,loigiai={Từ giả thiết $\overrightarrow{AM}=\dfrac{1}{3}\overrightarrow{AB}$ nên $M=A+\dfrac{1}{3}(B-A)=(\dfrac{14}{3};-1;\dfrac{7}{3})$.}]{Trong không gian $Oxyz$, cho $A(5;-3;3)$, $B(4;3;1)$. Tọa độ điểm $M$ thuộc đoạn thẳng $AB$ sao cho $2\overrightarrow{AM}=1\overrightarrow{MB}$ là}
  {$(\dfrac{14}{3};-1;\dfrac{7}{3})$.}
  {$(-\dfrac{14}{3};1;-\dfrac{7}{3})$.}{$(\dfrac{14}{3};\dfrac{7}{3};-1)$.}{$(\dfrac{11}{3};-1;\dfrac{7}{3})$.}

\cauTN[id=C2-B8-D09-TH-039,muc=TH,loigiai={Từ giả thiết $\overrightarrow{AM}=\dfrac{3}{4}\overrightarrow{AB}$ nên $M=A+\dfrac{3}{4}(B-A)=(2;-\dfrac{13}{4};-3)$.}]{Trong không gian $Oxyz$, cho $A(-4;-1;0)$, $B(4;-4;-4)$. Tọa độ điểm $M$ thuộc đoạn thẳng $AB$ sao cho $1\overrightarrow{AM}=3\overrightarrow{MB}$ là}
  {$(2;-\dfrac{13}{4};-3)$.}
  {$(-\dfrac{13}{4};2;-3)$.}{$(-2;-\dfrac{13}{4};-3)$.}{$(2;-\dfrac{13}{4};-5)$.}

\cauTN[id=C2-B8-D09-TH-040,muc=TH,loigiai={Từ giả thiết $\overrightarrow{AM}=2\overrightarrow{AB}$ nên $M=A+2(B-A)=(7;-2;3)$.}]{Trong không gian $Oxyz$, cho $A(-3;-4;-3)$, $B(2;-3;0)$. Tọa độ điểm $M$ thoả mãn $\overrightarrow{AM}=2\overrightarrow{AB}$ là}
  {$(7;-2;3)$.}
  {$(8;-2;3)$.}{$(5;-2;3)$.}{$(7;0;3)$.}

\cauTN[id=C2-B8-D09-TH-041,muc=TH,loigiai={Từ giả thiết $\overrightarrow{AM}=-1\overrightarrow{AB}$ nên $M=A-1(B-A)=(6;-7;-1)$.}]{Trong không gian $Oxyz$, cho $A(1;-5;0)$, $B(-4;-3;1)$. Tọa độ điểm $M$ thoả mãn $\overrightarrow{AM}=-1\overrightarrow{AB}$ là}
  {$(6;-7;-1)$.}
  {$(6;-7;1)$.}{$(-1;-7;6)$.}{$(5;-7;-1)$.}

\cauTN[id=C2-B8-D09-TH-042,muc=TH,loigiai={Từ giả thiết $\overrightarrow{AM}=3\overrightarrow{AB}$ nên $M=A+3(B-A)=(8;4;-5)$.}]{Trong không gian $Oxyz$, cho $A(-4;1;-5)$, $B(0;2;-5)$. Tọa độ điểm $M$ thoả mãn $\overrightarrow{AM}=3\overrightarrow{AB}$ là}
  {$(8;4;-5)$.}
  {$(10;4;-5)$.}{$(8;4;-7)$.}{$(8;4;-6)$.}

\cauTN[id=C2-B8-D09-TH-043,muc=TH,loigiai={Từ giả thiết $\overrightarrow{AM}=-1\overrightarrow{AB}$ nên $M=A-1(B-A)=(9;-3;1)$.}]{Trong không gian $Oxyz$, cho $A(4;1;2)$, $B(-1;5;3)$. Tọa độ điểm $M$ thoả mãn $\overrightarrow{AM}=-1\overrightarrow{AB}$ là}
  {$(9;-3;1)$.}
  {$(9;3;1)$.}{$(8;-3;1)$.}{$(10;-3;1)$.}

\cauTN[id=C2-B8-D09-TH-044,muc=TH,loigiai={Từ giả thiết $\overrightarrow{AM}=\dfrac{3}{4}\overrightarrow{AB}$ nên $M=A+\dfrac{3}{4}(B-A)=(-\dfrac{3}{2};-\dfrac{1}{4};-\dfrac{7}{4})$.}]{Trong không gian $Oxyz$, cho $A(0;-4;-1)$, $B(-2;1;-2)$. Tọa độ điểm $M$ thuộc đoạn thẳng $AB$ sao cho $1\overrightarrow{AM}=3\overrightarrow{MB}$ là}
  {$(-\dfrac{3}{2};-\dfrac{1}{4};-\dfrac{7}{4})$.}
  {$(-\dfrac{1}{2};-\dfrac{11}{4};-\dfrac{5}{4})$.}{$(-\dfrac{3}{2};-\dfrac{1}{4};\dfrac{7}{4})$.}{$(-\dfrac{3}{2};-\dfrac{1}{4};-\dfrac{11}{4})$.}

\cauTN[id=C2-B8-D09-TH-045,muc=TH,loigiai={Từ giả thiết $\overrightarrow{AM}=\dfrac{1}{4}\overrightarrow{AB}$ nên $M=A+\dfrac{1}{4}(B-A)=(\dfrac{3}{2};\dfrac{3}{2};-\dfrac{15}{4})$.}]{Trong không gian $Oxyz$, cho $A(3;3;-5)$, $B(-3;-3;0)$. Tọa độ điểm $M$ thuộc đoạn thẳng $AB$ sao cho $3\overrightarrow{AM}=1\overrightarrow{MB}$ là}
  {$(\dfrac{3}{2};\dfrac{3}{2};-\dfrac{15}{4})$.}
  {$(-\dfrac{3}{2};-\dfrac{3}{2};-\dfrac{5}{4})$.}{$(\dfrac{3}{2};\dfrac{7}{2};-\dfrac{15}{4})$.}{$(\dfrac{5}{2};\dfrac{3}{2};-\dfrac{15}{4})$.}

\cauTN[id=C2-B8-D09-TH-046,muc=TH,loigiai={Từ giả thiết $\overrightarrow{AM}=\dfrac{2}{3}\overrightarrow{AB}$ nên $M=A+\dfrac{2}{3}(B-A)=(2;\dfrac{8}{3};4)$.}]{Trong không gian $Oxyz$, cho $A(2;-2;2)$, $B(2;5;5)$. Tọa độ điểm $M$ thuộc đoạn thẳng $AB$ sao cho $1\overrightarrow{AM}=2\overrightarrow{MB}$ là}
  {$(2;\dfrac{8}{3};4)$.}
  {$(2;\dfrac{8}{3};5)$.}{$(-2;-\dfrac{8}{3};-4)$.}{$(1;\dfrac{8}{3};4)$.}

\cauTN[id=C2-B8-D09-TH-047,muc=TH,loigiai={Từ giả thiết $\overrightarrow{AM}=\dfrac{3}{4}\overrightarrow{AB}$ nên $M=A+\dfrac{3}{4}(B-A)=(0;-\dfrac{5}{4};2)$.}]{Trong không gian $Oxyz$, cho $A(0;4;5)$, $B(0;-3;1)$. Tọa độ điểm $M$ thuộc đoạn thẳng $AB$ sao cho $1\overrightarrow{AM}=3\overrightarrow{MB}$ là}
  {$(0;-\dfrac{5}{4};2)$.}
  {$(0;\dfrac{5}{4};-2)$.}{$(0;-\dfrac{9}{4};2)$.}{$(0;-\dfrac{5}{4};3)$.}

\cauTN[id=C2-B8-D09-TH-048,muc=TH,loigiai={Từ giả thiết $\overrightarrow{AM}=3\overrightarrow{AB}$ nên $M=A+3(B-A)=(-10;9;8)$.}]{Trong không gian $Oxyz$, cho $A(5;3;2)$, $B(0;5;4)$. Tọa độ điểm $M$ thoả mãn $\overrightarrow{AM}=3\overrightarrow{AB}$ là}
  {$(-10;9;8)$.}
  {$(-10;9;7)$.}{$(-11;9;8)$.}{$(-10;9;10)$.}

\cauDS[id=C2-B8-D09-TH-049,muc=TH,loigiai={Kiểm tra từng mệnh đề bằng công thức tọa độ. a) Sai. b) Sai. c) Sai. d) Đúng.}]{Trong không gian $Oxyz$, cho $A(-1;0;-1)$, $B(-2;-3;3)$. Xét tính đúng sai của các mệnh đề sau:}{%
  \y{Tọa độ điểm $M$ thoả mãn $\overrightarrow{AM}=-2\overrightarrow{AB}$ là $(3;6;-9)$.}{s}
  \y{Hình bình hành $ABCD$ có $A(-1;0;-1)$, $B(-2;-3;3)$, $C(2;-1;-2)$ thì $BD=\sqrt{130}$.}{s}
  \y{Tọa độ điểm $M$ thuộc đoạn thẳng $AB$ sao cho $3\overrightarrow{AM}=1\overrightarrow{MB}$ là $(-5;-12;15)$.}{s}
  \y{Với $A(2;4;3)$, $B(-3;-3;4)$, $C(7;11;2)$ thẳng hàng thì $\overrightarrow{AC}=-1\,\overrightarrow{AB}$.}{d}}

\cauDS[id=C2-B8-D09-TH-050,muc=TH,loigiai={Kiểm tra từng mệnh đề bằng công thức tọa độ. a) Sai. b) Sai. c) Đúng. d) Sai.}]{Trong không gian $Oxyz$, cho $A(4;-3;5)$, $B(-4;3;1)$, $C(1;2;-1)$. Xét tính đúng sai của các mệnh đề sau:}{%
  \y{Hình bình hành $ABCD$ có $A(4;-3;5)$, $B(-4;3;1)$, $C(1;2;-1)$ thì $BD=\sqrt{226}$.}{s}
  \y{Điểm $D$ để $ABCD$ là hình bình hành có tọa độ $(7;-4;3)$.}{s}
  \y{Tọa độ điểm $M$ thuộc đoạn thẳng $AB$ sao cho $3\overrightarrow{AM}=1\overrightarrow{MB}$ là $(2;-\dfrac{3}{2};4)$.}{d}
  \y{Hình bình hành $ABCD$ có $A(4;-3;5)$, $B(-4;3;1)$, $C(1;2;-1)$ thì $BD=\sqrt{221}$.}{s}}

\cauDS[id=C2-B8-D09-TH-051,muc=TH,loigiai={Kiểm tra từng mệnh đề bằng công thức tọa độ. a) Đúng. b) Đúng. c) Sai. d) Sai.}]{Trong không gian $Oxyz$, cho $A(0;4;-1)$, $B(5;-3;-3)$, $C(0;-4;-2)$. Xét tính đúng sai của các mệnh đề sau:}{%
  \y{Hình bình hành $ABCD$ có $A(0;4;-1)$, $B(5;-3;-3)$, $C(0;-4;-2)$ thì $BD=\sqrt{145}$.}{d}
  \y{Tọa độ điểm $M$ thuộc đoạn thẳng $AB$ sao cho $1\overrightarrow{AM}=3\overrightarrow{MB}$ là $(\dfrac{15}{4};-\dfrac{5}{4};-\dfrac{5}{2})$.}{d}
  \y{Tọa độ điểm $M$ thuộc đoạn thẳng $AB$ sao cho $3\overrightarrow{AM}=1\overrightarrow{MB}$ là $(-\dfrac{3}{4};\dfrac{9}{4};-\dfrac{3}{2})$.}{s}
  \y{Điểm $D$ để $ABCD$ là hình bình hành có tọa độ $(5;5;-2)$.}{s}}

\cauDS[id=C2-B8-D09-TH-052,muc=TH,loigiai={Kiểm tra từng mệnh đề bằng công thức tọa độ. a) Đúng. b) Sai. c) Đúng. d) Sai.}]{Trong không gian $Oxyz$, cho $A(2;2;-2)$, $B(-1;2;2)$, $C(-2;-1;1)$. Xét tính đúng sai của các mệnh đề sau:}{%
  \y{Hình bình hành $ABCD$ có $A(2;2;-2)$, $B(-1;2;2)$, $C(-2;-1;1)$ thì $BD=\sqrt{38}$.}{d}
  \y{Tọa độ điểm $M$ thoả mãn $\overrightarrow{AM}=2\overrightarrow{AB}$ là $(-4;0;6)$.}{s}
  \y{Điểm $D$ để $ABCD$ là hình bình hành có tọa độ $(1;-1;-3)$.}{d}
  \y{Với $A(-4;3;0)$, $B(-4;1;-2)$, $C(-4;0;-3)$ thẳng hàng thì $\overrightarrow{AC}=-\dfrac{3}{2}\,\overrightarrow{AB}$.}{s}}

\cauDS[id=C2-B8-D09-TH-053,muc=TH,loigiai={Kiểm tra từng mệnh đề bằng công thức tọa độ. a) Đúng. b) Sai. c) Sai. d) Sai.}]{Trong không gian $Oxyz$, cho $A(-3;1;-5)$, $B(4;-4;4)$. Xét tính đúng sai của các mệnh đề sau:}{%
  \y{Với $A(1;-2;-2)$, $B(-1;-3;-4)$, $C(3;-1;0)$ thẳng hàng thì $\overrightarrow{AC}=-1\,\overrightarrow{AB}$.}{d}
  \y{Hình bình hành $ABCD$ có $A(-3;1;-5)$, $B(4;-4;4)$, $C(5;-5;4)$ thì $BD=\sqrt{134}$.}{s}
  \y{Hình bình hành $ABCD$ có $A(-3;1;-5)$, $B(4;-4;4)$, $C(5;-5;4)$ thì $BD=\sqrt{137}$.}{s}
  \y{Tọa độ điểm $M$ thuộc đoạn thẳng $AB$ sao cho $1\overrightarrow{AM}=3\overrightarrow{MB}$ là $(\dfrac{9}{4};-\dfrac{11}{4};\dfrac{15}{4})$.}{s}}

\cauDS[id=C2-B8-D09-TH-054,muc=TH,loigiai={Kiểm tra từng mệnh đề bằng công thức tọa độ. a) Sai. b) Sai. c) Sai. d) Đúng.}]{Trong không gian $Oxyz$, cho $A(0;2;-1)$, $B(-4;-1;-1)$, $C(-2;2;3)$. Xét tính đúng sai của các mệnh đề sau:}{%
  \y{Với $A(-3;-3;1)$, $B(-1;3;-2)$, $C(-5;-9;4)$ thẳng hàng thì $\overrightarrow{AC}=1\,\overrightarrow{AB}$.}{s}
  \y{Tọa độ điểm $M$ thuộc đoạn thẳng $AB$ sao cho $2\overrightarrow{AM}=1\overrightarrow{MB}$ là $(-1;1;-\dfrac{4}{3})$.}{s}
  \y{Tọa độ điểm $M$ thoả mãn $\overrightarrow{AM}=3\overrightarrow{AB}$ là $(-14;-7;-1)$.}{s}
  \y{Điểm $D$ để $ABCD$ là hình bình hành có tọa độ $(2;5;3)$.}{d}}

\cauDS[id=C2-B8-D09-TH-055,muc=TH,loigiai={Kiểm tra từng mệnh đề bằng công thức tọa độ. a) Đúng. b) Đúng. c) Đúng. d) Sai.}]{Trong không gian $Oxyz$, cho $A(-4;4;2)$, $B(-3;-5;-3)$, $C(0;5;-5)$. Xét tính đúng sai của các mệnh đề sau:}{%
  \y{Tọa độ điểm $M$ thuộc đoạn thẳng $AB$ sao cho $2\overrightarrow{AM}=1\overrightarrow{MB}$ là $(-\dfrac{11}{3};1;\dfrac{1}{3})$.}{d}
  \y{Với $A(-2;-1;3)$, $B(4;-2;0)$, $C(10;-3;-3)$ thẳng hàng thì $\overrightarrow{AC}=2\,\overrightarrow{AB}$.}{d}
  \y{Điểm $D$ để $ABCD$ là hình bình hành có tọa độ $(-1;14;0)$.}{d}
  \y{Tọa độ điểm $M$ thoả mãn $\overrightarrow{AM}=2\overrightarrow{AB}$ là $(-2;-15;-8)$.}{s}}

\cauDS[id=C2-B8-D09-TH-056,muc=TH,loigiai={Kiểm tra từng mệnh đề bằng công thức tọa độ. a) Sai. b) Sai. c) Đúng. d) Sai.}]{Trong không gian $Oxyz$, cho $A(1;1;1)$, $B(2;3;-5)$. Xét tính đúng sai của các mệnh đề sau:}{%
  \y{Hình bình hành $ABCD$ có $A(1;1;1)$, $B(2;3;-5)$, $C(-5;4;-5)$ thì $BD=10$.}{s}
  \y{Với $A(0;-1;4)$, $B(-1;-1;-3)$, $C(-3;-1;-17)$ thẳng hàng thì $\overrightarrow{AC}=9\,\overrightarrow{AB}$.}{s}
  \y{Tọa độ điểm $M$ thuộc đoạn thẳng $AB$ sao cho $1\overrightarrow{AM}=2\overrightarrow{MB}$ là $(\dfrac{5}{3};\dfrac{7}{3};-3)$.}{d}
  \y{Với $A(-2;-2;2)$, $B(-4;0;0)$, $C(-5;1;-1)$ thẳng hàng thì $\overrightarrow{AC}=\dfrac{1}{2}\,\overrightarrow{AB}$.}{s}}

\cauDS[id=C2-B8-D09-TH-057,muc=TH,loigiai={Kiểm tra từng mệnh đề bằng công thức tọa độ. a) Sai. b) Đúng. c) Sai. d) Đúng.}]{Trong không gian $Oxyz$, cho $A(-3;3;4)$, $B(-3;4;-4)$, $C(-5;-4;5)$. Xét tính đúng sai của các mệnh đề sau:}{%
  \y{Với $A(-3;-2;1)$, $B(1;3;-2)$, $C(5;8;-5)$ thẳng hàng thì $\overrightarrow{AC}=5\,\overrightarrow{AB}$.}{s}
  \y{Hình bình hành $ABCD$ có $A(-3;3;4)$, $B(-3;4;-4)$, $C(-5;-4;5)$ thì $BD=\sqrt{374}$.}{d}
  \y{Hình bình hành $ABCD$ có $A(-3;3;4)$, $B(-3;4;-4)$, $C(-5;-4;5)$ thì $BD=5\sqrt{15}$.}{s}
  \y{Điểm $D$ để $ABCD$ là hình bình hành có tọa độ $(-5;-5;13)$.}{d}}

\cauDS[id=C2-B8-D09-TH-058,muc=TH,loigiai={Kiểm tra từng mệnh đề bằng công thức tọa độ. a) Đúng. b) Đúng. c) Sai. d) Sai.}]{Trong không gian $Oxyz$, cho $A(-3;-3;0)$, $B(-2;1;4)$, $C(-4;1;3)$. Xét tính đúng sai của các mệnh đề sau:}{%
  \y{Điểm $D$ để $ABCD$ là hình bình hành có tọa độ $(-5;-3;-1)$.}{d}
  \y{Với $A(1;-3;2)$, $B(3;2;3)$, $C(7;12;5)$ thẳng hàng thì $\overrightarrow{AC}=3\,\overrightarrow{AB}$.}{d}
  \y{Hình bình hành $ABCD$ có $A(-3;-3;0)$, $B(-2;1;4)$, $C(-4;1;3)$ thì $BD=3\sqrt{6}$.}{s}
  \y{Với $A(-1;3;-3)$, $B(1;1;1)$, $C(5;-3;9)$ thẳng hàng thì $\overrightarrow{AC}=\dfrac{6004799503160661}{18014398509481984}\,\overrightarrow{AB}$.}{s}}

\cauDS[id=C2-B8-D09-TH-059,muc=TH,loigiai={Kiểm tra từng mệnh đề bằng công thức tọa độ. a) Đúng. b) Đúng. c) Sai. d) Sai.}]{Trong không gian $Oxyz$, cho $A(-2;2;0)$, $B(-2;4;-2)$. Xét tính đúng sai của các mệnh đề sau:}{%
  \y{Tọa độ điểm $M$ thuộc đoạn thẳng $AB$ sao cho $1\overrightarrow{AM}=2\overrightarrow{MB}$ là $(-2;\dfrac{10}{3};-\dfrac{4}{3})$.}{d}
  \y{Hình bình hành $ABCD$ có $A(-2;2;0)$, $B(-2;4;-2)$, $C(-2;-1;4)$ thì $BD=\sqrt{113}$.}{d}
  \y{Với $A(1;2;-4)$, $B(-1;3;0)$, $C(3;1;-8)$ thẳng hàng thì $\overrightarrow{AC}=1\,\overrightarrow{AB}$.}{s}
  \y{Hình bình hành $ABCD$ có $A(-2;2;0)$, $B(-2;4;-2)$, $C(-2;-1;4)$ thì $BD=4\sqrt{7}$.}{s}}

\cauDS[id=C2-B8-D09-TH-060,muc=TH,loigiai={Kiểm tra từng mệnh đề bằng công thức tọa độ. a) Sai. b) Đúng. c) Đúng. d) Đúng.}]{Trong không gian $Oxyz$, cho $A(3;1;-4)$, $B(4;0;4)$, $C(-2;1;1)$. Xét tính đúng sai của các mệnh đề sau:}{%
  \y{Điểm $D$ để $ABCD$ là hình bình hành có tọa độ $(9;0;-1)$.}{s}
  \y{Với $A(-4;-4;1)$, $B(-2;-1;-1)$, $C(-8;-10;5)$ thẳng hàng thì $\overrightarrow{AC}=-2\,\overrightarrow{AB}$.}{d}
  \y{Hình bình hành $ABCD$ có $A(3;1;-4)$, $B(4;0;4)$, $C(-2;1;1)$ thì $BD=\sqrt{174}$.}{d}
  \y{Với $A(-1;2;3)$, $B(-2;2;-4)$, $C(0;2;10)$ thẳng hàng thì $\overrightarrow{AC}=-1\,\overrightarrow{AB}$.}{d}}

\cauDS[id=C2-B8-D09-TH-061,muc=TH,loigiai={Kiểm tra từng mệnh đề bằng công thức tọa độ. a) Đúng. b) Sai. c) Sai. d) Sai.}]{Trong không gian $Oxyz$, cho $A(5;3;4)$, $B(4;-5;-2)$, $C(0;-1;5)$. Xét tính đúng sai của các mệnh đề sau:}{%
  \y{Hình bình hành $ABCD$ có $A(5;3;4)$, $B(4;-5;-2)$, $C(0;-1;5)$ thì $BD=\sqrt{322}$.}{d}
  \y{Điểm $D$ để $ABCD$ là hình bình hành có tọa độ $(-1;-9;-1)$.}{s}
  \y{Hình bình hành $ABCD$ có $A(5;3;4)$, $B(4;-5;-2)$, $C(0;-1;5)$ thì $BD=\sqrt{321}$.}{s}
  \y{Với $A(0;0;-4)$, $B(-4;-4;4)$, $C(-2;-2;0)$ thẳng hàng thì $\overrightarrow{AC}=-\dfrac{1}{2}\,\overrightarrow{AB}$.}{s}}

\cauDS[id=C2-B8-D09-TH-062,muc=TH,loigiai={Kiểm tra từng mệnh đề bằng công thức tọa độ. a) Đúng. b) Đúng. c) Đúng. d) Sai.}]{Trong không gian $Oxyz$, cho $A(3;-2;2)$, $B(-5;-2;-5)$, $C(4;-4;1)$. Xét tính đúng sai của các mệnh đề sau:}{%
  \y{Điểm $D$ để $ABCD$ là hình bình hành có tọa độ $(12;-4;8)$.}{d}
  \y{Với $A(-2;-1;-2)$, $B(1;3;-3)$, $C(-5;-5;-1)$ thẳng hàng thì $\overrightarrow{AC}=-1\,\overrightarrow{AB}$.}{d}
  \y{Với $A(-4;4;-1)$, $B(2;2;4)$, $C(-16;8;-11)$ thẳng hàng thì $\overrightarrow{AC}=-2\,\overrightarrow{AB}$.}{d}
  \y{Tọa độ điểm $M$ thuộc đoạn thẳng $AB$ sao cho $1\overrightarrow{AM}=3\overrightarrow{MB}$ là $(-3;-2;-\dfrac{5}{4})$.}{s}}

\cauDS[id=C2-B8-D09-TH-063,muc=TH,loigiai={Kiểm tra từng mệnh đề bằng công thức tọa độ. a) Đúng. b) Sai. c) Đúng. d) Đúng.}]{Trong không gian $Oxyz$, cho $A(0;3;-5)$, $B(-2;-2;4)$, $C(-1;-3;5)$. Xét tính đúng sai của các mệnh đề sau:}{%
  \y{Điểm $D$ để $ABCD$ là hình bình hành có tọa độ $(1;2;-4)$.}{d}
  \y{Hình bình hành $ABCD$ có $A(0;3;-5)$, $B(-2;-2;4)$, $C(-1;-3;5)$ thì $BD=2\sqrt{22}$.}{s}
  \y{Với $A(-2;0;2)$, $B(3;4;3)$, $C(-7;-4;1)$ thẳng hàng thì $\overrightarrow{AC}=-1\,\overrightarrow{AB}$.}{d}
  \y{Hình bình hành $ABCD$ có $A(0;3;-5)$, $B(-2;-2;4)$, $C(-1;-3;5)$ thì $BD=\sqrt{89}$.}{d}}

\cauDS[id=C2-B8-D09-TH-064,muc=TH,loigiai={Kiểm tra từng mệnh đề bằng công thức tọa độ. a) Đúng. b) Đúng. c) Đúng. d) Sai.}]{Trong không gian $Oxyz$, cho $A(-1;0;-4)$, $B(3;-5;2)$, $C(2;-1;-2)$. Xét tính đúng sai của các mệnh đề sau:}{%
  \y{Với $A(-1;-4;-1)$, $B(2;-2;2)$, $C(-4;-6;-4)$ thẳng hàng thì $\overrightarrow{AC}=-1\,\overrightarrow{AB}$.}{d}
  \y{Hình bình hành $ABCD$ có $A(-1;0;-4)$, $B(3;-5;2)$, $C(2;-1;-2)$ thì $BD=\sqrt{206}$.}{d}
  \y{Điểm $D$ để $ABCD$ là hình bình hành có tọa độ $(-2;4;-8)$.}{d}
  \y{Hình bình hành $ABCD$ có $A(-1;0;-4)$, $B(3;-5;2)$, $C(2;-1;-2)$ thì $BD=\sqrt{210}$.}{s}}

\cauTLN[id=C2-B8-D09-TH-065,muc=TH,loigiai={Ta có $\overrightarrow{AB}=(-4;1;7)$, $\overrightarrow{AC}=(-8;2;14)$. Vì $\overrightarrow{AC}=k\overrightarrow{AB}$ nên $k=2$.}]{Cho ba điểm thẳng hàng $A(0;2;-3)$, $B(-4;3;4)$, $C(-8;4;11)$. Biết $\overrightarrow{AC}=k\overrightarrow{AB}$. Tính $k$.}{2}

\cauTLN[id=C2-B8-D09-TH-066,muc=TH,loigiai={Ta có $\overrightarrow{AB}=(-2;-4;0)$, $\overrightarrow{AC}=(4;8;0)$. Vì $\overrightarrow{AC}=k\overrightarrow{AB}$ nên $k=-2$.}]{Cho ba điểm thẳng hàng $A(0;3;-2)$, $B(-2;-1;-2)$, $C(4;11;-2)$. Biết $\overrightarrow{AC}=k\overrightarrow{AB}$. Tính $k$.}{-2}

\cauTLN[id=C2-B8-D09-TH-067,muc=TH,loigiai={Ta có $\overrightarrow{AB}=(-7;-3;-6)$, $\overrightarrow{AC}=(-14;-6;-12)$. Vì $\overrightarrow{AC}=k\overrightarrow{AB}$ nên $k=2$.}]{Cho ba điểm thẳng hàng $A(3;-1;3)$, $B(-4;-4;-3)$, $C(-11;-7;-9)$. Biết $\overrightarrow{AC}=k\overrightarrow{AB}$. Tính $k$.}{2}

\cauTLN[id=C2-B8-D09-TH-068,muc=TH,loigiai={Ta có $\overrightarrow{AB}=(-7;2;1)$, $\overrightarrow{AC}=(-21;6;3)$. Vì $\overrightarrow{AC}=k\overrightarrow{AB}$ nên $k=3$.}]{Cho ba điểm thẳng hàng $A(3;-3;0)$, $B(-4;-1;1)$, $C(-18;3;3)$. Biết $\overrightarrow{AC}=k\overrightarrow{AB}$. Tính $k$.}{3}

\cauTLN[id=C2-B8-D09-TH-069,muc=TH,loigiai={Ta có $\overrightarrow{AB}=(1;1;-5)$, $\overrightarrow{AC}=(-1;-1;5)$. Vì $\overrightarrow{AC}=k\overrightarrow{AB}$ nên $k=-1$.}]{Cho ba điểm thẳng hàng $A(0;-2;1)$, $B(1;-1;-4)$, $C(-1;-3;6)$. Biết $\overrightarrow{AC}=k\overrightarrow{AB}$. Tính $k$.}{-1}

\cauTLN[id=C2-B8-D09-TH-070,muc=TH,loigiai={Ta có $\overrightarrow{AB}=(4;5;0)$, $\overrightarrow{AC}=(8;10;0)$. Vì $\overrightarrow{AC}=k\overrightarrow{AB}$ nên $k=2$.}]{Cho ba điểm thẳng hàng $A(-2;-1;2)$, $B(2;4;2)$, $C(6;9;2)$. Biết $\overrightarrow{AC}=k\overrightarrow{AB}$. Tính $k$.}{2}

\cauTLN[id=C2-B8-D09-TH-071,muc=TH,loigiai={Hình bình hành $ABCD$ có $D=A+C-B=(-1;-11;-8)$. Do đó $\overrightarrow{BD}=(-2;-14;-10)$ và $BD=10\sqrt{3}$.}]{Cho hình bình hành $ABCD$ với $A(-3;-5;-2)$, $B(1;3;2)$, $C(3;-3;-4)$. Tính $BD^2$.}{300}

\cauTLN[id=C2-B8-D09-TH-072,muc=TH,loigiai={Hình bình hành $ABCD$ có $D=A+C-B=(1;-8;-9)$. Do đó $\overrightarrow{BD}=(3;-12;-13)$ và $BD=\sqrt{322}$.}]{Cho hình bình hành $ABCD$ với $A(3;-3;-3)$, $B(-2;4;4)$, $C(-4;-1;-2)$. Tính $BD^2$.}{322}

\cauTLN[id=C2-B8-D09-TH-073,muc=TH,loigiai={Hình bình hành $ABCD$ có $D=A+C-B=(-9;-1;3)$. Do đó $\overrightarrow{BD}=(-14;-3;1)$ và $BD=\sqrt{206}$.}]{Cho hình bình hành $ABCD$ với $A(-4;-2;4)$, $B(5;2;2)$, $C(0;3;1)$. Tính $BD^2$.}{206}

\cauTLN[id=C2-B8-D09-TH-074,muc=TH,loigiai={Hình bình hành $ABCD$ có $D=A+C-B=(4;0;-9)$. Do đó $\overrightarrow{BD}=(9;0;-8)$ và $BD=\sqrt{145}$.}]{Cho hình bình hành $ABCD$ với $A(0;3;-5)$, $B(-5;0;-1)$, $C(-1;-3;-5)$. Tính $BD^2$.}{145}

\cauTLN[id=C2-B8-D09-TH-075,muc=TH,loigiai={Hình bình hành $ABCD$ có $D=A+C-B=(-1;-7;10)$. Do đó $\overrightarrow{BD}=(-3;-7;15)$ và $BD=\sqrt{283}$.}]{Cho hình bình hành $ABCD$ với $A(-2;-4;0)$, $B(2;0;-5)$, $C(3;-3;5)$. Tính $BD^2$.}{283}

\cauTLN[id=C2-B8-D09-TH-076,muc=TH,loigiai={Ba điểm $A,B,C$ thẳng hàng nên $\overrightarrow{AC}=k\overrightarrow{AB}$. Từ hoành độ: $k=-2$, suy ra $z_C=-1-2\cdot(1)=-3$.}]{Cho ba điểm $A(4;-2;-1)$, $B(-3;-4;0)$, $C(18;2;c)$ thẳng hàng. Tính $c$.}{-3}

\cauTLN[id=C2-B8-D09-TH-077,muc=TH,loigiai={Ba điểm $A,B,C$ thẳng hàng nên $\overrightarrow{AC}=k\overrightarrow{AB}$. Từ hoành độ: $k=-1$, suy ra $z_C=-2-1\cdot(2)=-4$.}]{Cho ba điểm $A(4;-2;-2)$, $B(-4;-2;0)$, $C(12;-2;c)$ thẳng hàng. Tính $c$.}{-4}

\cauTLN[id=C2-B8-D09-TH-078,muc=TH,loigiai={Ba điểm $A,B,C$ thẳng hàng nên $\overrightarrow{AC}=k\overrightarrow{AB}$. Từ hoành độ: $k=3$, suy ra $z_C=3+3\cdot(-7)=-18$.}]{Cho ba điểm $A(4;1;3)$, $B(-2;-4;-4)$, $C(-14;-14;c)$ thẳng hàng. Tính $c$.}{-18}

\cauTLN[id=C2-B8-D09-TH-079,muc=TH,loigiai={Ba điểm $A,B,C$ thẳng hàng nên $\overrightarrow{AC}=k\overrightarrow{AB}$. Từ hoành độ: $k=-1$, suy ra $z_C=-2-1\cdot(6)=-8$.}]{Cho ba điểm $A(1;-4;-2)$, $B(2;3;4)$, $C(0;-11;c)$ thẳng hàng. Tính $c$.}{-8}

\cauTLN[id=C2-B8-D09-TH-080,muc=TH,loigiai={Ba điểm $A,B,C$ thẳng hàng nên $\overrightarrow{AC}=k\overrightarrow{AB}$. Từ hoành độ: $k=2$, suy ra $z_C=-4+2\cdot(3)=2$.}]{Cho ba điểm $A(2;-4;-4)$, $B(0;-4;-1)$, $C(-2;-4;c)$ thẳng hàng. Tính $c$.}{2}

\cauTN[id=C2-B8-D09-VD-001,muc=VD,loigiai={Gọi $M(0;0;t)$. Ba điểm $A,B,M$ thẳng hàng nên $\overrightarrow{AM}$ cùng phương $\overrightarrow{AB}$. Giải ra $t=2$, tức là $M(0;0;2)$.}]{Cho hai điểm $A(6;-2;-2)$, $B(-3;1;4)$. Điểm $M$ thuộc trục $Oz$ sao cho $A$, $B$, $M$ thẳng hàng có tọa độ là}
  {$(0;0;2)$.}
  {$(0;1;2)$.}{$(0;0;1)$.}{$(0;0;4)$.}

\cauTN[id=C2-B8-D09-VD-002,muc=VD,loigiai={Gọi $M(0;0;t)$. Ba điểm $A,B,M$ thẳng hàng nên $\overrightarrow{AM}$ cùng phương $\overrightarrow{AB}$. Giải ra $t=3$, tức là $M(0;0;3)$.}]{Cho hai điểm $A(-2;6;7)$, $B(2;-6;-1)$. Điểm $M$ thuộc trục $Oz$ sao cho $A$, $B$, $M$ thẳng hàng có tọa độ là}
  {$(0;0;3)$.}
  {$(0;0;4)$.}{$(-2;0;3)$.}{$(-1;0;3)$.}

\cauTN[id=C2-B8-D09-VD-003,muc=VD,loigiai={Gọi $M(0;0;t)$. Ba điểm $A,B,M$ thẳng hàng nên $\overrightarrow{AM}$ cùng phương $\overrightarrow{AB}$. Giải ra $t=0$, tức là $M(0;0;0)$.}]{Cho hai điểm $A(3;-9;-9)$, $B(1;-3;-3)$. Điểm $M$ thuộc trục $Oz$ sao cho $A$, $B$, $M$ thẳng hàng có tọa độ là}
  {$(0;0;0)$.}
  {$(1;0;0)$.}{$(0;0;-1)$.}{$(-2;0;0)$.}

\cauTN[id=C2-B8-D09-VD-004,muc=VD,loigiai={Gọi $M(0;0;t)$. Ba điểm $A,B,M$ thẳng hàng nên $\overrightarrow{AM}$ cùng phương $\overrightarrow{AB}$. Giải ra $t=-2$, tức là $M(0;0;-2)$.}]{Cho hai điểm $A(2;2;1)$, $B(-4;-4;-8)$. Điểm $M$ thuộc trục $Oz$ sao cho $A$, $B$, $M$ thẳng hàng có tọa độ là}
  {$(0;0;-2)$.}
  {$(0;0;2)$.}{$(-1;0;-2)$.}{$(0;1;-2)$.}

\cauTN[id=C2-B8-D09-VD-005,muc=VD,loigiai={Gọi $M(0;0;t)$. Ba điểm $A,B,M$ thẳng hàng nên $\overrightarrow{AM}$ cùng phương $\overrightarrow{AB}$. Giải ra $t=-4$, tức là $M(0;0;-4)$.}]{Cho hai điểm $A(3;-1;-1)$, $B(6;-2;2)$. Điểm $M$ thuộc trục $Oz$ sao cho $A$, $B$, $M$ thẳng hàng có tọa độ là}
  {$(0;0;-4)$.}
  {$(0;-4;0)$.}{$(-1;0;-4)$.}{$(-2;0;-4)$.}

\cauTN[id=C2-B8-D09-VD-006,muc=VD,loigiai={Gọi $M(0;0;t)$. Ba điểm $A,B,M$ thẳng hàng nên $\overrightarrow{AM}$ cùng phương $\overrightarrow{AB}$. Giải ra $t=-3$, tức là $M(0;0;-3)$.}]{Cho hai điểm $A(3;2;-1)$, $B(6;4;1)$. Điểm $M$ thuộc trục $Oz$ sao cho $A$, $B$, $M$ thẳng hàng có tọa độ là}
  {$(0;0;-3)$.}
  {$(0;1;-3)$.}{$(0;0;-1)$.}{$(0;-3;0)$.}

\cauTN[id=C2-B8-D09-VD-007,muc=VD,loigiai={Gọi $M(0;0;t)$. Ba điểm $A,B,M$ thẳng hàng nên $\overrightarrow{AM}$ cùng phương $\overrightarrow{AB}$. Giải ra $t=2$, tức là $M(0;0;2)$.}]{Cho hai điểm $A(-6;-9;11)$, $B(2;3;-1)$. Điểm $M$ thuộc trục $Oz$ sao cho $A$, $B$, $M$ thẳng hàng có tọa độ là}
  {$(0;0;2)$.}
  {$(0;-1;2)$.}{$(0;0;3)$.}{$(0;0;0)$.}

\cauTN[id=C2-B8-D09-VD-008,muc=VD,loigiai={Gọi $M(0;0;t)$. Ba điểm $A,B,M$ thẳng hàng nên $\overrightarrow{AM}$ cùng phương $\overrightarrow{AB}$. Giải ra $t=1$, tức là $M(0;0;1)$.}]{Cho hai điểm $A(-9;9;-8)$, $B(3;-3;4)$. Điểm $M$ thuộc trục $Oz$ sao cho $A$, $B$, $M$ thẳng hàng có tọa độ là}
  {$(0;0;1)$.}
  {$(2;0;1)$.}{$(0;0;-1)$.}{$(0;0;3)$.}

\cauTN[id=C2-B8-D09-VD-009,muc=VD,loigiai={Gọi $M(0;0;t)$. Ba điểm $A,B,M$ thẳng hàng nên $\overrightarrow{AM}$ cùng phương $\overrightarrow{AB}$. Giải ra $t=1$, tức là $M(0;0;1)$.}]{Cho hai điểm $A(-2;-2;1)$, $B(-3;-3;1)$. Điểm $M$ thuộc trục $Oz$ sao cho $A$, $B$, $M$ thẳng hàng có tọa độ là}
  {$(0;0;1)$.}
  {$(0;0;-1)$.}{$(0;0;3)$.}{$(0;1;1)$.}

\cauTN[id=C2-B8-D09-VD-010,muc=VD,loigiai={Gọi $M(0;0;t)$. Ba điểm $A,B,M$ thẳng hàng nên $\overrightarrow{AM}$ cùng phương $\overrightarrow{AB}$. Giải ra $t=4$, tức là $M(0;0;4)$.}]{Cho hai điểm $A(-1;-2;2)$, $B(2;4;8)$. Điểm $M$ thuộc trục $Oz$ sao cho $A$, $B$, $M$ thẳng hàng có tọa độ là}
  {$(0;0;4)$.}
  {$(0;2;4)$.}{$(-1;0;4)$.}{$(0;0;6)$.}

\cauTN[id=C2-B8-D09-VD-011,muc=VD,loigiai={Gọi $M(0;0;t)$. Ba điểm $A,B,M$ thẳng hàng nên $\overrightarrow{AM}$ cùng phương $\overrightarrow{AB}$. Giải ra $t=-3$, tức là $M(0;0;-3)$.}]{Cho hai điểm $A(-2;6;-7)$, $B(3;-9;3)$. Điểm $M$ thuộc trục $Oz$ sao cho $A$, $B$, $M$ thẳng hàng có tọa độ là}
  {$(0;0;-3)$.}
  {$(-3;0;0)$.}{$(2;0;-3)$.}{$(0;0;-5)$.}

\cauTN[id=C2-B8-D09-VD-012,muc=VD,loigiai={Gọi $M(0;0;t)$. Ba điểm $A,B,M$ thẳng hàng nên $\overrightarrow{AM}$ cùng phương $\overrightarrow{AB}$. Giải ra $t=-3$, tức là $M(0;0;-3)$.}]{Cho hai điểm $A(-2;-2;1)$, $B(1;1;-5)$. Điểm $M$ thuộc trục $Oz$ sao cho $A$, $B$, $M$ thẳng hàng có tọa độ là}
  {$(0;0;-3)$.}
  {$(0;0;-5)$.}{$(0;-1;-3)$.}{$(0;0;3)$.}

\cauTN[id=C2-B8-D09-VD-013,muc=VD,loigiai={Hình bình hành $ABCD$ có $D=A+C-B=(-5;5;-11)$. Do đó $\overrightarrow{BD}=(-9;9;-16)$ và $BD=\sqrt{418}$.}]{Cho hình bình hành $ABCD$ với $A(-5;0;0)$, $B(4;-4;5)$, $C(4;1;-6)$. Độ dài đường chéo $BD$ bằng}
  {$\sqrt{418}$.}
  {$\sqrt{419}$.}{$\sqrt{417}$.}{$\sqrt{422}$.}

\cauTN[id=C2-B8-D09-VD-014,muc=VD,loigiai={Hình bình hành $ABCD$ có $D=A+C-B=(-1;4;6)$. Do đó $\overrightarrow{BD}=(-5;2;9)$ và $BD=\sqrt{110}$.}]{Cho hình bình hành $ABCD$ với $A(-1;0;5)$, $B(4;2;-3)$, $C(4;6;-2)$. Độ dài đường chéo $BD$ bằng}
  {$\sqrt{110}$.}
  {$\sqrt{111}$.}{$\sqrt{109}$.}{$\sqrt{114}$.}

\cauTN[id=C2-B8-D09-VD-015,muc=VD,loigiai={Hình bình hành $ABCD$ có $D=A+C-B=(-3;-1;-2)$. Do đó $\overrightarrow{BD}=(-5;0;-6)$ và $BD=\sqrt{61}$.}]{Cho hình bình hành $ABCD$ với $A(-2;2;3)$, $B(2;-1;4)$, $C(1;-4;-1)$. Độ dài đường chéo $BD$ bằng}
  {$\sqrt{61}$.}
  {$\sqrt{62}$.}{$2\sqrt{15}$.}{$\sqrt{65}$.}

\cauTN[id=C2-B8-D09-VD-016,muc=VD,loigiai={Hình bình hành $ABCD$ có $D=A+C-B=(-2;-3;-2)$. Do đó $\overrightarrow{BD}=(2;-5;4)$ và $BD=3\sqrt{5}$.}]{Cho hình bình hành $ABCD$ với $A(0;0;-6)$, $B(-4;2;-6)$, $C(-6;-1;-2)$. Độ dài đường chéo $BD$ bằng}
  {$3\sqrt{5}$.}
  {$\sqrt{46}$.}{$2\sqrt{11}$.}{$7$.}

\cauTN[id=C2-B8-D09-VD-017,muc=VD,loigiai={Hình bình hành $ABCD$ có $D=A+C-B=(-3;1;-2)$. Do đó $\overrightarrow{BD}=(-6;-1;-4)$ và $BD=\sqrt{53}$.}]{Cho hình bình hành $ABCD$ với $A(3;-1;-4)$, $B(3;2;2)$, $C(-3;4;4)$. Độ dài đường chéo $BD$ bằng}
  {$\sqrt{53}$.}
  {$3\sqrt{6}$.}{$2\sqrt{13}$.}{$\sqrt{57}$.}

\cauTN[id=C2-B8-D09-VD-018,muc=VD,loigiai={Hình bình hành $ABCD$ có $D=A+C-B=(8;-9;4)$. Do đó $\overrightarrow{BD}=(13;-15;6)$ và $BD=\sqrt{430}$.}]{Cho hình bình hành $ABCD$ với $A(0;-1;1)$, $B(-5;6;-2)$, $C(3;-2;1)$. Độ dài đường chéo $BD$ bằng}
  {$\sqrt{430}$.}
  {$\sqrt{431}$.}{$\sqrt{429}$.}{$\sqrt{434}$.}

\cauTN[id=C2-B8-D09-VD-019,muc=VD,loigiai={Hình bình hành $ABCD$ có $D=A+C-B=(-1;7;2)$. Do đó $\overrightarrow{BD}=(-1;9;4)$ và $BD=7\sqrt{2}$.}]{Cho hình bình hành $ABCD$ với $A(-2;1;-3)$, $B(0;-2;-2)$, $C(1;4;3)$. Độ dài đường chéo $BD$ bằng}
  {$7\sqrt{2}$.}
  {$3\sqrt{11}$.}{$\sqrt{97}$.}{$\sqrt{102}$.}

\cauTN[id=C2-B8-D09-VD-020,muc=VD,loigiai={Hình bình hành $ABCD$ có $D=A+C-B=(4;4;0)$. Do đó $\overrightarrow{BD}=(10;6;-2)$ và $BD=2\sqrt{35}$.}]{Cho hình bình hành $ABCD$ với $A(-6;2;5)$, $B(-6;-2;2)$, $C(4;0;-3)$. Độ dài đường chéo $BD$ bằng}
  {$2\sqrt{35}$.}
  {$\sqrt{141}$.}{$\sqrt{139}$.}{$12$.}

\cauTN[id=C2-B8-D09-VD-021,muc=VD,loigiai={Hình bình hành $ABCD$ có $D=A+C-B=(9;15;-2)$. Do đó $\overrightarrow{BD}=(15;18;-3)$ và $BD=3\sqrt{62}$.}]{Cho hình bình hành $ABCD$ với $A(6;6;4)$, $B(-6;-3;1)$, $C(-3;6;-5)$. Độ dài đường chéo $BD$ bằng}
  {$3\sqrt{62}$.}
  {$\sqrt{559}$.}{$\sqrt{557}$.}{$\sqrt{562}$.}

\cauTN[id=C2-B8-D09-VD-022,muc=VD,loigiai={Hình bình hành $ABCD$ có $D=A+C-B=(-6;13;7)$. Do đó $\overrightarrow{BD}=(-10;17;7)$ và $BD=\sqrt{438}$.}]{Cho hình bình hành $ABCD$ với $A(2;4;6)$, $B(4;-4;0)$, $C(-4;5;1)$. Độ dài đường chéo $BD$ bằng}
  {$\sqrt{438}$.}
  {$\sqrt{439}$.}{$\sqrt{437}$.}{$\sqrt{442}$.}

\cauTN[id=C2-B8-D09-VD-023,muc=VD,loigiai={Hình bình hành $ABCD$ có $D=A+C-B=(-2;-1;-11)$. Do đó $\overrightarrow{BD}=(1;-5;-17)$ và $BD=3\sqrt{35}$.}]{Cho hình bình hành $ABCD$ với $A(1;3;-1)$, $B(-3;4;6)$, $C(-6;0;-4)$. Độ dài đường chéo $BD$ bằng}
  {$3\sqrt{35}$.}
  {$2\sqrt{79}$.}{$\sqrt{314}$.}{$\sqrt{319}$.}

\cauTN[id=C2-B8-D09-VD-024,muc=VD,loigiai={Hình bình hành $ABCD$ có $D=A+C-B=(0;-11;-12)$. Do đó $\overrightarrow{BD}=(-6;-17;-16)$ và $BD=\sqrt{581}$.}]{Cho hình bình hành $ABCD$ với $A(0;-3;-3)$, $B(6;6;4)$, $C(6;-2;-5)$. Độ dài đường chéo $BD$ bằng}
  {$\sqrt{581}$.}
  {$\sqrt{582}$.}{$2\sqrt{145}$.}{$3\sqrt{65}$.}

\cauTN[id=C2-B8-D09-VD-025,muc=VD,loigiai={Ta có $\overrightarrow{AB}=(6;-1;-6)$, $\overrightarrow{AC}=(-6;1;6)$. Vì $\overrightarrow{AC}=k\overrightarrow{AB}$ nên $k=-1$.}]{Cho ba điểm thẳng hàng $A(-4;-3;4)$, $B(2;-4;-2)$, $C(-10;-2;10)$. Biết $\overrightarrow{AC}=k\overrightarrow{AB}$. Giá trị của $k$ bằng}
  {$-1$.}
  {$1$.}{$2$.}{$-3$.}

\cauTN[id=C2-B8-D09-VD-026,muc=VD,loigiai={Ta có $\overrightarrow{AB}=(-1;-2;5)$, $\overrightarrow{AC}=(-3;-6;15)$. Vì $\overrightarrow{AC}=k\overrightarrow{AB}$ nên $k=3$.}]{Cho ba điểm thẳng hàng $A(-1;3;-4)$, $B(-2;1;1)$, $C(-4;-3;11)$. Biết $\overrightarrow{AC}=k\overrightarrow{AB}$. Giá trị của $k$ bằng}
  {$3$.}
  {$-2$.}{$5$.}{$6$.}

\cauTN[id=C2-B8-D09-VD-027,muc=VD,loigiai={Ta có $\overrightarrow{AB}=(-6;4;-2)$, $\overrightarrow{AC}=(3;-2;1)$. Vì $\overrightarrow{AC}=k\overrightarrow{AB}$ nên $k=-\dfrac{1}{2}$.}]{Cho ba điểm thẳng hàng $A(3;-3;4)$, $B(-3;1;2)$, $C(6;-5;5)$. Biết $\overrightarrow{AC}=k\overrightarrow{AB}$. Giá trị của $k$ bằng}
  {$-\dfrac{1}{2}$.}
  {$-\dfrac{1}{4}$.}{$-2$.}{$-1$.}

\cauTN[id=C2-B8-D09-VD-028,muc=VD,loigiai={Ta có $\overrightarrow{AB}=(4;3;1)$, $\overrightarrow{AC}=(12;9;3)$. Vì $\overrightarrow{AC}=k\overrightarrow{AB}$ nên $k=3$.}]{Cho ba điểm thẳng hàng $A(-2;-3;-2)$, $B(2;0;-1)$, $C(10;6;1)$. Biết $\overrightarrow{AC}=k\overrightarrow{AB}$. Giá trị của $k$ bằng}
  {$3$.}
  {$9$.}{$6$.}{$-3$.}

\cauTN[id=C2-B8-D09-VD-029,muc=VD,loigiai={Ta có $\overrightarrow{AB}=(-7;0;-2)$, $\overrightarrow{AC}=(7;0;2)$. Vì $\overrightarrow{AC}=k\overrightarrow{AB}$ nên $k=-1$.}]{Cho ba điểm thẳng hàng $A(3;-4;-2)$, $B(-4;-4;-4)$, $C(10;-4;0)$. Biết $\overrightarrow{AC}=k\overrightarrow{AB}$. Giá trị của $k$ bằng}
  {$-1$.}
  {$1$.}{$-4$.}{$2$.}

\cauTN[id=C2-B8-D09-VD-030,muc=VD,loigiai={Ta có $\overrightarrow{AB}=(2;-4;0)$, $\overrightarrow{AC}=(1;-2;0)$. Vì $\overrightarrow{AC}=k\overrightarrow{AB}$ nên $k=\dfrac{1}{2}$.}]{Cho ba điểm thẳng hàng $A(0;4;-4)$, $B(2;0;-4)$, $C(1;2;-4)$. Biết $\overrightarrow{AC}=k\overrightarrow{AB}$. Giá trị của $k$ bằng}
  {$\dfrac{1}{2}$.}
  {$-\dfrac{5}{2}$.}{$2$.}{$\dfrac{5}{2}$.}

\cauTN[id=C2-B8-D09-VD-031,muc=VD,loigiai={Ta có $\overrightarrow{AB}=(-5;8;-1)$, $\overrightarrow{AC}=(-15;24;-3)$. Vì $\overrightarrow{AC}=k\overrightarrow{AB}$ nên $k=3$.}]{Cho ba điểm thẳng hàng $A(3;-4;4)$, $B(-2;4;3)$, $C(-12;20;1)$. Biết $\overrightarrow{AC}=k\overrightarrow{AB}$. Giá trị của $k$ bằng}
  {$3$.}
  {$4$.}{$2$.}{$-2$.}

\cauTN[id=C2-B8-D09-VD-032,muc=VD,loigiai={Ta có $\overrightarrow{AB}=(4;2;-1)$, $\overrightarrow{AC}=(12;6;-3)$. Vì $\overrightarrow{AC}=k\overrightarrow{AB}$ nên $k=3$.}]{Cho ba điểm thẳng hàng $A(-1;1;1)$, $B(3;3;0)$, $C(11;7;-2)$. Biết $\overrightarrow{AC}=k\overrightarrow{AB}$. Giá trị của $k$ bằng}
  {$3$.}
  {$\dfrac{6004799503160661}{18014398509481984}$.}{$-3$.}{$4$.}

\cauTN[id=C2-B8-D09-VD-033,muc=VD,loigiai={Ta có $\overrightarrow{AB}=(-4;4;2)$, $\overrightarrow{AC}=(8;-8;-4)$. Vì $\overrightarrow{AC}=k\overrightarrow{AB}$ nên $k=-2$.}]{Cho ba điểm thẳng hàng $A(3;-3;-4)$, $B(-1;1;-2)$, $C(11;-11;-8)$. Biết $\overrightarrow{AC}=k\overrightarrow{AB}$. Giá trị của $k$ bằng}
  {$-2$.}
  {$-1$.}{$-3$.}{$3$.}

\cauTN[id=C2-B8-D09-VD-034,muc=VD,loigiai={Ta có $\overrightarrow{AB}=(2;-2;1)$, $\overrightarrow{AC}=(-4;4;-2)$. Vì $\overrightarrow{AC}=k\overrightarrow{AB}$ nên $k=-2$.}]{Cho ba điểm thẳng hàng $A(0;4;1)$, $B(2;2;2)$, $C(-4;8;-1)$. Biết $\overrightarrow{AC}=k\overrightarrow{AB}$. Giá trị của $k$ bằng}
  {$-2$.}
  {$-3$.}{$2$.}{$-4$.}

\cauTN[id=C2-B8-D09-VD-035,muc=VD,loigiai={Ta có $\overrightarrow{AB}=(1;7;-7)$, $\overrightarrow{AC}=(2;14;-14)$. Vì $\overrightarrow{AC}=k\overrightarrow{AB}$ nên $k=2$.}]{Cho ba điểm thẳng hàng $A(1;-3;4)$, $B(2;4;-3)$, $C(3;11;-10)$. Biết $\overrightarrow{AC}=k\overrightarrow{AB}$. Giá trị của $k$ bằng}
  {$2$.}
  {$\dfrac{1}{2}$.}{$5$.}{$-2$.}

\cauTN[id=C2-B8-D09-VD-036,muc=VD,loigiai={Ta có $\overrightarrow{AB}=(4;1;1)$, $\overrightarrow{AC}=(12;3;3)$. Vì $\overrightarrow{AC}=k\overrightarrow{AB}$ nên $k=3$.}]{Cho ba điểm thẳng hàng $A(-1;2;1)$, $B(3;3;2)$, $C(11;5;4)$. Biết $\overrightarrow{AC}=k\overrightarrow{AB}$. Giá trị của $k$ bằng}
  {$3$.}
  {$9$.}{$\dfrac{3}{2}$.}{$-3$.}

\cauDS[id=C2-B8-D09-VD-037,muc=VD,loigiai={Kiểm tra từng mệnh đề bằng công thức tọa độ. a) Đúng. b) Sai. c) Đúng. d) Đúng.}]{Trong không gian $Oxyz$, cho $A(0;-5;6)$, $B(1;4;2)$. Xét tính đúng sai của các mệnh đề sau:}{%
  \y{Tọa độ điểm $M$ thoả mãn $\overrightarrow{AM}=2\overrightarrow{AB}$ là $(2;13;-2)$.}{d}
  \y{Với $A(2;3;4)$, $B(4;-4;-2)$, $C(-2;17;16)$ thẳng hàng thì $\overrightarrow{AC}=-4\,\overrightarrow{AB}$.}{s}
  \y{Hình bình hành $ABCD$ có $A(0;-5;6)$, $B(1;4;2)$, $C(5;4;-6)$ thì $BD=\sqrt{106}$.}{d}
  \y{Với $A(-2;0;4)$, $B(-2;2;4)$, $C(-2;1;4)$ thẳng hàng thì $\overrightarrow{AC}=\dfrac{1}{2}\,\overrightarrow{AB}$.}{d}}

\cauDS[id=C2-B8-D09-VD-038,muc=VD,loigiai={Kiểm tra từng mệnh đề bằng công thức tọa độ. a) Sai. b) Đúng. c) Sai. d) Sai.}]{Trong không gian $Oxyz$, cho $A(-1;-2;5)$, $B(-2;4;4)$, $C(-6;-3;1)$. Xét tính đúng sai của các mệnh đề sau:}{%
  \y{Với $A(3;-2;2)$, $B(-1;-4;-1)$, $C(11;2;8)$ thẳng hàng thì $\overrightarrow{AC}=3\,\overrightarrow{AB}$.}{s}
  \y{Hình bình hành $ABCD$ có $A(-1;-2;5)$, $B(-2;4;4)$, $C(-6;-3;1)$ thì $BD=\sqrt{182}$.}{d}
  \y{Điểm $D$ để $ABCD$ là hình bình hành có tọa độ $(-5;-9;0)$.}{s}
  \y{Điểm $D$ để $ABCD$ là hình bình hành có tọa độ $(-5;-11;2)$.}{s}}

\cauDS[id=C2-B8-D09-VD-039,muc=VD,loigiai={Kiểm tra từng mệnh đề bằng công thức tọa độ. a) Sai. b) Đúng. c) Sai. d) Đúng.}]{Trong không gian $Oxyz$, cho $A(1;5;2)$, $B(0;3;-3)$, $C(1;0;2)$. Xét tính đúng sai của các mệnh đề sau:}{%
  \y{Hình bình hành $ABCD$ có $A(1;5;2)$, $B(0;3;-3)$, $C(1;0;2)$ thì $BD=\sqrt{109}$.}{s}
  \y{Tọa độ điểm $M$ thuộc đoạn thẳng $AB$ sao cho $3\overrightarrow{AM}=1\overrightarrow{MB}$ là $(\dfrac{3}{4};\dfrac{9}{2};\dfrac{3}{4})$.}{d}
  \y{Với $A(-4;-1;3)$, $B(4;-2;3)$, $C(20;-4;3)$ thẳng hàng thì $\overrightarrow{AC}=\dfrac{3}{2}\,\overrightarrow{AB}$.}{s}
  \y{Điểm $D$ để $ABCD$ là hình bình hành có tọa độ $(2;2;7)$.}{d}}

\cauDS[id=C2-B8-D09-VD-040,muc=VD,loigiai={Kiểm tra từng mệnh đề bằng công thức tọa độ. a) Đúng. b) Sai. c) Sai. d) Đúng.}]{Trong không gian $Oxyz$, cho $A(0;-4;2)$, $B(-4;2;0)$, $C(-3;4;0)$. Xét tính đúng sai của các mệnh đề sau:}{%
  \y{Điểm $D$ để $ABCD$ là hình bình hành có tọa độ $(1;-2;2)$.}{d}
  \y{Hình bình hành $ABCD$ có $A(0;-4;2)$, $B(-4;2;0)$, $C(-3;4;0)$ thì $BD=2\sqrt{11}$.}{s}
  \y{Với $A(2;-3;-1)$, $B(1;-3;3)$, $C(4;-3;-9)$ thẳng hàng thì $\overrightarrow{AC}=1\,\overrightarrow{AB}$.}{s}
  \y{Với $A(0;-4;0)$, $B(2;3;-1)$, $C(6;17;-3)$ thẳng hàng thì $\overrightarrow{AC}=3\,\overrightarrow{AB}$.}{d}}

\cauDS[id=C2-B8-D09-VD-041,muc=VD,loigiai={Kiểm tra từng mệnh đề bằng công thức tọa độ. a) Sai. b) Sai. c) Đúng. d) Đúng.}]{Trong không gian $Oxyz$, cho $A(5;6;-6)$, $B(2;-4;6)$, $C(-6;5;-1)$. Xét tính đúng sai của các mệnh đề sau:}{%
  \y{Tọa độ điểm $M$ thoả mãn $\overrightarrow{AM}=-2\overrightarrow{AB}$ là $(11;-26;-30)$.}{s}
  \y{Hình bình hành $ABCD$ có $A(5;6;-6)$, $B(2;-4;6)$, $C(-6;5;-1)$ thì $BD=\sqrt{751}$.}{s}
  \y{Với $A(0;1;4)$, $B(2;-4;-4)$, $C(6;-14;-20)$ thẳng hàng thì $\overrightarrow{AC}=3\,\overrightarrow{AB}$.}{d}
  \y{Điểm $D$ để $ABCD$ là hình bình hành có tọa độ $(-3;15;-13)$.}{d}}

\cauDS[id=C2-B8-D09-VD-042,muc=VD,loigiai={Kiểm tra từng mệnh đề bằng công thức tọa độ. a) Sai. b) Sai. c) Đúng. d) Sai.}]{Trong không gian $Oxyz$, cho $A(6;-5;-6)$, $B(2;-2;-4)$, $C(6;-4;-1)$. Xét tính đúng sai của các mệnh đề sau:}{%
  \y{Hình bình hành $ABCD$ có $A(6;-5;-6)$, $B(2;-2;-4)$, $C(6;-4;-1)$ thì $BD=\sqrt{94}$.}{s}
  \y{Hình bình hành $ABCD$ có $A(6;-5;-6)$, $B(2;-2;-4)$, $C(6;-4;-1)$ thì $BD=\sqrt{91}$.}{s}
  \y{Với $A(-4;1;2)$, $B(-4;-1;4)$, $C(-4;-3;6)$ thẳng hàng thì $\overrightarrow{AC}=2\,\overrightarrow{AB}$.}{d}
  \y{Điểm $D$ để $ABCD$ là hình bình hành có tọa độ $(10;-7;-2)$.}{s}}

\cauDS[id=C2-B8-D09-VD-043,muc=VD,loigiai={Kiểm tra từng mệnh đề bằng công thức tọa độ. a) Đúng. b) Đúng. c) Đúng. d) Sai.}]{Trong không gian $Oxyz$, cho $A(-4;-3;-3)$, $B(1;-3;2)$, $C(-6;1;1)$. Xét tính đúng sai của các mệnh đề sau:}{%
  \y{Với $A(-2;-3;2)$, $B(-3;3;2)$, $C(-4;9;2)$ thẳng hàng thì $\overrightarrow{AC}=2\,\overrightarrow{AB}$.}{d}
  \y{Tọa độ điểm $M$ thuộc đoạn thẳng $AB$ sao cho $3\overrightarrow{AM}=1\overrightarrow{MB}$ là $(-\dfrac{11}{4};-3;-\dfrac{7}{4})$.}{d}
  \y{Với $A(-1;2;2)$, $B(2;2;-1)$, $C(8;2;-7)$ thẳng hàng thì $\overrightarrow{AC}=3\,\overrightarrow{AB}$.}{d}
  \y{Điểm $D$ để $ABCD$ là hình bình hành có tọa độ $(11;-1;4)$.}{s}}

\cauDS[id=C2-B8-D09-VD-044,muc=VD,loigiai={Kiểm tra từng mệnh đề bằng công thức tọa độ. a) Sai. b) Sai. c) Sai. d) Đúng.}]{Trong không gian $Oxyz$, cho $A(1;-5;1)$, $B(-6;-1;4)$, $C(-6;0;0)$. Xét tính đúng sai của các mệnh đề sau:}{%
  \y{Tọa độ điểm $M$ thuộc đoạn thẳng $AB$ sao cho $1\overrightarrow{AM}=2\overrightarrow{MB}$ là $(-\dfrac{8}{3};-\dfrac{7}{3};3)$.}{s}
  \y{Tọa độ điểm $M$ thoả mãn $\overrightarrow{AM}=-1\overrightarrow{AB}$ là $(8;9;-2)$.}{s}
  \y{Điểm $D$ để $ABCD$ là hình bình hành có tọa độ $(-4;1;-3)$.}{s}
  \y{Với $A(-1;3;4)$, $B(-3;0;3)$, $C(3;9;6)$ thẳng hàng thì $\overrightarrow{AC}=-2\,\overrightarrow{AB}$.}{d}}

\cauDS[id=C2-B8-D09-VD-045,muc=VD,loigiai={Kiểm tra từng mệnh đề bằng công thức tọa độ. a) Sai. b) Đúng. c) Đúng. d) Sai.}]{Trong không gian $Oxyz$, cho $A(-3;5;6)$, $B(-3;3;-3)$, $C(-1;2;-3)$. Xét tính đúng sai của các mệnh đề sau:}{%
  \y{Điểm $D$ để $ABCD$ là hình bình hành có tọa độ $(-1;4;7)$.}{s}
  \y{Tọa độ điểm $M$ thoả mãn $\overrightarrow{AM}=-2\overrightarrow{AB}$ là $(-3;9;24)$.}{d}
  \y{Tọa độ điểm $M$ thuộc đoạn thẳng $AB$ sao cho $2\overrightarrow{AM}=1\overrightarrow{MB}$ là $(-3;\dfrac{13}{3};3)$.}{d}
  \y{Với $A(3;1;2)$, $B(-1;1;-1)$, $C(-9;1;-7)$ thẳng hàng thì $\overrightarrow{AC}=6\,\overrightarrow{AB}$.}{s}}

\cauDS[id=C2-B8-D09-VD-046,muc=VD,loigiai={Kiểm tra từng mệnh đề bằng công thức tọa độ. a) Đúng. b) Sai. c) Sai. d) Sai.}]{Trong không gian $Oxyz$, cho $A(-5;3;-6)$, $B(5;-6;-5)$, $C(0;1;0)$. Xét tính đúng sai của các mệnh đề sau:}{%
  \y{Tọa độ điểm $M$ thuộc đoạn thẳng $AB$ sao cho $3\overrightarrow{AM}=1\overrightarrow{MB}$ là $(-\dfrac{5}{2};\dfrac{3}{4};-\dfrac{23}{4})$.}{d}
  \y{Với $A(-3;3;-3)$, $B(4;0;0)$, $C(11;-3;3)$ thẳng hàng thì $\overrightarrow{AC}=4\,\overrightarrow{AB}$.}{s}
  \y{Điểm $D$ để $ABCD$ là hình bình hành có tọa độ $(-10;9;-1)$.}{s}
  \y{Tọa độ điểm $M$ thoả mãn $\overrightarrow{AM}=-2\overrightarrow{AB}$ là $(-25;23;-8)$.}{s}}

\cauDS[id=C2-B8-D09-VD-047,muc=VD,loigiai={Kiểm tra từng mệnh đề bằng công thức tọa độ. a) Đúng. b) Sai. c) Đúng. d) Đúng.}]{Trong không gian $Oxyz$, cho $A(3;-6;-4)$, $B(-3;0;-6)$, $C(6;-6;-4)$. Xét tính đúng sai của các mệnh đề sau:}{%
  \y{Với $A(0;-4;-4)$, $B(1;4;2)$, $C(2;12;8)$ thẳng hàng thì $\overrightarrow{AC}=2\,\overrightarrow{AB}$.}{d}
  \y{Hình bình hành $ABCD$ có $A(3;-6;-4)$, $B(-3;0;-6)$, $C(6;-6;-4)$ thì $BD=\sqrt{389}$.}{s}
  \y{Với $A(2;-3;-2)$, $B(2;2;-2)$, $C(2;-13;-2)$ thẳng hàng thì $\overrightarrow{AC}=-2\,\overrightarrow{AB}$.}{d}
  \y{Điểm $D$ để $ABCD$ là hình bình hành có tọa độ $(12;-12;-2)$.}{d}}

\cauDS[id=C2-B8-D09-VD-048,muc=VD,loigiai={Kiểm tra từng mệnh đề bằng công thức tọa độ. a) Sai. b) Sai. c) Sai. d) Đúng.}]{Trong không gian $Oxyz$, cho $A(-3;2;4)$, $B(3;3;-6)$, $C(-6;-5;5)$. Xét tính đúng sai của các mệnh đề sau:}{%
  \y{Hình bình hành $ABCD$ có $A(-3;2;4)$, $B(3;3;-6)$, $C(-6;-5;5)$ thì $BD=2\sqrt{187}$.}{s}
  \y{Tọa độ điểm $M$ thuộc đoạn thẳng $AB$ sao cho $1\overrightarrow{AM}=3\overrightarrow{MB}$ là $(\dfrac{3}{2};\dfrac{15}{4};-\dfrac{7}{2})$.}{s}
  \y{Điểm $D$ để $ABCD$ là hình bình hành có tọa độ $(-12;-4;15)$.}{s}
  \y{Điểm $D$ để $ABCD$ là hình bình hành có tọa độ $(-12;-6;15)$.}{d}}

\cauTLN[id=C2-B8-D09-VD-049,muc=VD,loigiai={Gọi $M(0;0;t)$. Ba điểm $A,B,M$ thẳng hàng nên $\overrightarrow{AM}$ cùng phương $\overrightarrow{AB}$. Giải ra $t=0$, tức là $M(0;0;0)$.}]{Cho hai điểm $A(3;-1;-2)$, $B(-3;1;2)$. Điểm $M$ thuộc trục $Oz$ sao cho $A,B,M$ thẳng hàng. Tính $z_M$.}{0}

\cauTLN[id=C2-B8-D09-VD-050,muc=VD,loigiai={Gọi $M(0;0;t)$. Ba điểm $A,B,M$ thẳng hàng nên $\overrightarrow{AM}$ cùng phương $\overrightarrow{AB}$. Giải ra $t=2$, tức là $M(0;0;2)$.}]{Cho hai điểm $A(6;-6;6)$, $B(-9;9;-4)$. Điểm $M$ thuộc trục $Oz$ sao cho $A,B,M$ thẳng hàng. Tính $z_M$.}{2}

\cauTLN[id=C2-B8-D09-VD-051,muc=VD,loigiai={Gọi $M(0;0;t)$. Ba điểm $A,B,M$ thẳng hàng nên $\overrightarrow{AM}$ cùng phương $\overrightarrow{AB}$. Giải ra $t=2$, tức là $M(0;0;2)$.}]{Cho hai điểm $A(1;1;4)$, $B(-2;-2;-2)$. Điểm $M$ thuộc trục $Oz$ sao cho $A,B,M$ thẳng hàng. Tính $z_M$.}{2}

\cauTLN[id=C2-B8-D09-VD-052,muc=VD,loigiai={Gọi $M(0;0;t)$. Ba điểm $A,B,M$ thẳng hàng nên $\overrightarrow{AM}$ cùng phương $\overrightarrow{AB}$. Giải ra $t=3$, tức là $M(0;0;3)$.}]{Cho hai điểm $A(-9;9;9)$, $B(-3;3;5)$. Điểm $M$ thuộc trục $Oz$ sao cho $A,B,M$ thẳng hàng. Tính $z_M$.}{3}

\cauTLN[id=C2-B8-D09-VD-053,muc=VD,loigiai={Gọi $M(0;0;t)$. Ba điểm $A,B,M$ thẳng hàng nên $\overrightarrow{AM}$ cùng phương $\overrightarrow{AB}$. Giải ra $t=-2$, tức là $M(0;0;-2)$.}]{Cho hai điểm $A(-3;2;-1)$, $B(-6;4;0)$. Điểm $M$ thuộc trục $Oz$ sao cho $A,B,M$ thẳng hàng. Tính $z_M$.}{-2}

\cauTLN[id=C2-B8-D09-VD-054,muc=VD,loigiai={Gọi $M(0;0;t)$. Ba điểm $A,B,M$ thẳng hàng nên $\overrightarrow{AM}$ cùng phương $\overrightarrow{AB}$. Giải ra $t=4$, tức là $M(0;0;4)$.}]{Cho hai điểm $A(-1;-3;4)$, $B(2;6;4)$. Điểm $M$ thuộc trục $Oz$ sao cho $A,B,M$ thẳng hàng. Tính $z_M$.}{4}

\cauTLN[id=C2-B8-D09-VD-055,muc=VD,loigiai={Hình bình hành $ABCD$ có $D=A+C-B=(9;12;4)$. Do đó $\overrightarrow{BD}=(13;17;8)$ và $BD=3\sqrt{58}$.}]{Cho hình bình hành $ABCD$ với $A(6;1;6)$, $B(-4;-5;-4)$, $C(-1;6;-6)$. Tính $BD^2$.}{522}

\cauTLN[id=C2-B8-D09-VD-056,muc=VD,loigiai={Hình bình hành $ABCD$ có $D=A+C-B=(-11;8;-4)$. Do đó $\overrightarrow{BD}=(-15;12;-8)$ và $BD=\sqrt{433}$.}]{Cho hình bình hành $ABCD$ với $A(-4;2;5)$, $B(4;-4;4)$, $C(-3;2;-5)$. Tính $BD^2$.}{433}

\cauTLN[id=C2-B8-D09-VD-057,muc=VD,loigiai={Hình bình hành $ABCD$ có $D=A+C-B=(6;4;-3)$. Do đó $\overrightarrow{BD}=(4;5;-6)$ và $BD=\sqrt{77}$.}]{Cho hình bình hành $ABCD$ với $A(4;2;2)$, $B(2;-1;3)$, $C(4;1;-2)$. Tính $BD^2$.}{77}

\cauTLN[id=C2-B8-D09-VD-058,muc=VD,loigiai={Hình bình hành $ABCD$ có $D=A+C-B=(1;1;0)$. Do đó $\overrightarrow{BD}=(7;-5;-1)$ và $BD=5\sqrt{3}$.}]{Cho hình bình hành $ABCD$ với $A(-6;1;3)$, $B(-6;6;1)$, $C(1;6;-2)$. Tính $BD^2$.}{75}

\cauTLN[id=C2-B8-D09-VD-059,muc=VD,loigiai={Hình bình hành $ABCD$ có $D=A+C-B=(-4;4;-1)$. Do đó $\overrightarrow{BD}=(-5;-2;-1)$ và $BD=\sqrt{30}$.}]{Cho hình bình hành $ABCD$ với $A(-4;4;-3)$, $B(1;6;0)$, $C(1;6;2)$. Tính $BD^2$.}{30}

\cauTLN[id=C2-B8-D09-VD-060,muc=VD,loigiai={Hình bình hành $ABCD$ có $D=A+C-B=(-4;-14;16)$. Do đó $\overrightarrow{BD}=(-7;-19;20)$ và $BD=9\sqrt{10}$.}]{Cho hình bình hành $ABCD$ với $A(-4;-6;6)$, $B(3;5;-4)$, $C(3;-3;6)$. Tính $BD^2$.}{810}
